\input{preamble}

% OK, start here.
%
\begin{document}

\title{Algèbre à puissances divisées}


\maketitle

\phantomsection
\label{section-phantom}

\tableofcontents

\section{Introduction}
\label{section-introduction}

\noindent
Dans ce chapitre, nous étudions les algèbres à puissances divisées et ce que
l'on peut en faire. L'ouvrage \cite{Berthelot} constitue une référence.





\section{Puissances divisées}
\label{section-divided-powers}

\noindent
Nous rassemblons dans cette section quelques résultats sur les anneaux à
puissances divisées. Nous adoptons la convention $0! = 1$ (puisqu'un produit
vide doit valoir $1$).

\begin{definition}
\label{definition-divided-powers}
Soit $A$ un anneau. Soit $I$ un idéal de $A$. Une famille d'applications
$\gamma_n : I \to I$, $n > 0$, est appelée {\it structure de puissances divisées}
sur $I$ si, pour tous $n \geq 0$, $m > 0$, $x, y \in I$ et $a \in A$, on a
\begin{enumerate}
\item $\gamma_1(x) = x$, et l'on pose également $\gamma_0(x) = 1$,
\item $\gamma_n(x)\gamma_m(x) = \frac{(n + m)!}{n! m!} \gamma_{n + m}(x)$,
\item $\gamma_n(ax) = a^n \gamma_n(x)$,
\item $\gamma_n(x + y) = \sum_{i = 0, \ldots, n} \gamma_i(x)\gamma_{n - i}(y)$,
\item $\gamma_n(\gamma_m(x)) = \frac{(nm)!}{n! (m!)^n} \gamma_{nm}(x)$.
\end{enumerate}
\end{definition}

\noindent
Remarquons que les nombres rationnels $\frac{(n + m)!}{n! m!}$
et $\frac{(nm)!}{n! (m!)^n}$ qui interviennent dans la définition sont en
réalité des entiers : le premier est le nombre de façons de choisir $n$
éléments parmi $n + m$, et le second compte les façons de répartir $nm$
objets en $n$ groupes de $m$ objets.
Voici quelques remarques sur la définition qui montrent que
$\gamma_n(x)$ remplace $x^n/n!$ dans $I$.

\begin{lemma}
\label{lemma-silly}
Soit $A$ un anneau. Soit $I$ un idéal de $A$.
\begin{enumerate}
\item Si $\gamma$ est une structure de puissances divisées\footnote{Ici
et dans la suite, $\gamma$ désigne une suite d'applications
$\gamma_1, \gamma_2, \gamma_3, \ldots$ de $I$ dans $I$.}
sur $I$, alors
$n! \gamma_n(x) = x^n$ pour $n \geq 1$, $x \in I$.
\end{enumerate}
Supposons que $A$ soit sans torsion en tant que $\mathbf{Z}$-module.
\begin{enumerate}
\item[(2)] Une structure de puissances divisées sur $I$, si elle existe, est unique.
\item[(3)] Si $\gamma_n : I \to I$ sont des applications, alors
$$
\gamma\text{ est une structure de puissances divisées}
\Leftrightarrow
n! \gamma_n(x) = x^n\ \forall x \in I, n \geq 1.
$$
\item[(4)] L'idéal $I$ admet une structure de puissances divisées
si et seulement s'il existe
des éléments $x_i$ qui engendrent $I$ comme idéal et tels que,
pour tout $n \geq 1$, on ait $x_i^n \in (n!)I$.
\end{enumerate}
\end{lemma}

\begin{proof}
Démonstration de (1). Si $\gamma$ est une structure de puissances divisées,
alors la condition (2) (appliquée à $1$ et $n-1$ à la place de $n$ et $m$)
entraîne que $n \gamma_n(x) = \gamma_1(x)\gamma_{n - 1}(x)$. Par récurrence
et grâce à la condition (1), on obtient donc $n! \gamma_n(x) = x^n$.

\medskip\noindent
Supposons que $A$ soit sans torsion en tant que $\mathbf{Z}$-module.
Démonstration de (2). Cela résulte immédiatement de (1).

\medskip\noindent
Démonstration de (3). Supposons que $n! \gamma_n(x) = x^n$ pour tout
$x \in I$ et tout $n \geq 1$. Comme
$A \subset A \otimes_{\mathbf{Z}} \mathbf{Q}$, il suffit de démontrer les
axiomes (1) -- (5) de la définition \ref{definition-divided-powers} lorsque
$A$ est une $\mathbf{Q}$-algèbre. Dans ce cas,
$\gamma_n(x) = x^n/n!$, et la vérification de (1) -- (5) est immédiate;
par exemple, (4) correspond à la formule du binôme
$$
(x + y)^n = \sum_{i = 0, \ldots, n} \frac{n!}{i!(n - i)!} x^iy^{n - i}
$$
Nous laissons au lecteur le soin d'effectuer ces vérifications afin de
s'assurer que les coefficients sont corrects.

\medskip\noindent
Démonstration de (4). Supposons que des éléments $x_i$ engendrent $I$
comme idéal et vérifient $x_i^n \in (n!)I$ pour tout $n \geq 1$. Nous affirmons
que, pour tout $x \in I$, on a $x^n \in (n!)I$. Si cette assertion est vraie, on
peut poser $\gamma_n(x) = x^n/n!$, ce qui définit une structure de puissances
divisées d'après (3). Pour démontrer l'assertion, remarquons qu'elle est vraie
pour $x = ax_i$. Elle est donc vraie pour un ensemble d'éléments qui engendre
$I$ comme groupe abélien. Par récurrence sur la longueur d'une expression en
fonction de ces générateurs, il suffit de la vérifier pour $x + y$ lorsqu'elle
est vraie pour $x$ et $y$. Cela découle immédiatement de la formule du binôme.
\end{proof}

\begin{example}
\label{example-ideal-generated-by-p}
Soit $p$ un nombre premier.
Soit $A$ un anneau dans lequel tout entier $n$ non divisible par $p$
est inversible, autrement dit, $A$ est une $\mathbf{Z}_{(p)}$-algèbre. Alors
$I = pA$ admet une structure canonique de puissances divisées. Plus précisément,
pour $x = pa \in I$, on pose
$$
\gamma_n(x) = \frac{p^n}{n!} a^n
$$
On vérifie immédiatement que $p^n/n! \in p\mathbf{Z}_{(p)}$
pour $n \geq 1$ (cela découle par exemple du fait que l'exposant de $p$
dans la décomposition en facteurs premiers de $n!$ est
$\left\lfloor n/p \right\rfloor + \left\lfloor n/p^2 \right\rfloor
+ \left\lfloor n/p^3 \right\rfloor + \ldots$),
de sorte que la définition a un sens et fournit une suite d'applications
$\gamma_n : I \to I$. Il est immédiat de vérifier que les conditions
(1) -- (5) de la définition \ref{definition-divided-powers} sont satisfaites.
On peut aussi observer que la définition convient manifestement à
$A_0 = \mathbf{Z}_{(p)}$, puis appliquer le
lemme \ref{lemma-gamma-extends}.
\end{example}

\noindent
On remarque que $\gamma_n\left(0\right) = 0$ pour tout idéal $I$ de
$A$ et toute structure de puissances divisées $\gamma$ sur $I$. (Cela résulte
de l'axiome (3) de la définition \ref{definition-divided-powers},
en prenant $a=0$.)

\begin{lemma}
\label{lemma-check-on-generators}
Soit $A$ un anneau. Soit $I$ un idéal de $A$. Soit $\gamma_n : I \to I$,
$n \geq 1$, une suite d'applications. Supposons que
\begin{enumerate}
\item[(a)] les propriétés (1), (3) et (4) de la définition
\ref{definition-divided-powers} soient vérifiées pour tous $x, y \in I$, et
\item[(b)] les propriétés (2) et (5) soient vérifiées pour tout $x$ appartenant
à un ensemble d'éléments qui engendre $I$ comme idéal.
\end{enumerate}
Alors $\gamma$ est une structure de puissances divisées sur $I$.
\end{lemma}

\begin{proof}
Dans cette démonstration, les numéros (1), (2), (3), (4), (5) renvoient aux
conditions énoncées dans la définition \ref{definition-divided-powers}.
Il résulte de (3) que, si (2) et (5) sont vérifiées pour $x$, elles le sont
pour $ax$, quel que soit $a \in A$. L'hypothèse (b) entraîne donc que
(2) et (5) sont vérifiées pour un ensemble d'éléments qui engendre $I$
comme groupe abélien. Par récurrence sur la longueur d'une expression en
fonction de ces générateurs, il suffit par conséquent de montrer que, si
$x, y \in I$ et si (2) et (5) sont vérifiées pour $x$ et $y$, alors elles
le sont pour $x + y$.

\medskip\noindent
Démonstration de (2) pour $x + y$. D'après (4), on a
$$
\gamma_n(x + y)\gamma_m(x + y) =
\sum\nolimits_{i + j = n,\ k + l = m}
\gamma_i(x)\gamma_k(x)\gamma_j(y)\gamma_l(y)
$$
En utilisant (2) pour $x$ et $y$, ceci vaut
$$
\sum \frac{(i + k)!}{i!k!}\frac{(j + l)!}{j!l!}
\gamma_{i + k}(x)\gamma_{j + l}(y)
$$
En comparant avec le développement
$$
\gamma_{n + m}(x + y) = \sum \gamma_a(x)\gamma_b(y)
$$
on voit qu'il faut démontrer que, pour $a + b = n + m$, on a
$$
\sum\nolimits_{i + k = a,\ j + l = b,\ i + j = n,\ k + l = m}
\frac{(i + k)!}{i!k!}\frac{(j + l)!}{j!l!}
=
\frac{(n + m)!}{n!m!}.
$$
Plutôt que de démontrer directement cette identité, remarquons qu'elle est
vraie pour l'idéal $I = (x, y)$ de l'anneau de polynômes
$\mathbf{Q}[x, y]$, car $\gamma_n(f) = f^n/n!$, $f \in I$, définit une
structure de puissances divisées sur $I$. L'identité entre nombres rationnels
ci-dessus est donc vraie.

\medskip\noindent
Démonstration de (5) pour $x + y$, sachant que (1) -- (4) sont vérifiées et
que (5) l'est pour $x$ et $y$. Nous nous ramenons de nouveau à une identité
entre nombres rationnels. En effet, (4) permet d'écrire
$\gamma_n(\gamma_m(x + y)) = \gamma_n(\sum \gamma_i(x)\gamma_j(y))$.
Une nouvelle application de (4) permet d'écrire
$\gamma_n(\gamma_m(x + y))$ comme une somme de termes qui sont des produits
de facteurs de la forme $\gamma_k(\gamma_i(x)\gamma_j(y))$.
Si $i > 0$, alors
\begin{align*}
\gamma_k(\gamma_i(x)\gamma_j(y)) & =
\gamma_j(y)^k\gamma_k(\gamma_i(x)) \\
& = \frac{(ki)!}{k!(i!)^k} \gamma_j(y)^k \gamma_{ki}(x) \\
& =
\frac{(ki)!}{k!(i!)^k} \frac{(kj)!}{(j!)^k} \gamma_{ki}(x) \gamma_{kj}(y)
\end{align*}
où l'on a utilisé (3) dans la première égalité, (5) pour $x$ dans la
deuxième et (2) exactement $k$ fois dans la troisième. En utilisant (5)
pour $y$, on voit que la même égalité vaut lorsque $i = 0$. En poursuivant
ainsi et en utilisant tous les axiomes sauf (5), on obtient
$$
\gamma_n(\gamma_m(x + y)) =
\sum\nolimits_{i + j = nm} c_{ij}\gamma_i(x)\gamma_j(y)
$$
où les $c_{ij} \in \mathbf{Z}$ sont certaines constantes universelles. Comme
précédemment, le fait que l'égalité soit valable dans l'anneau de polynômes
$\mathbf{Q}[x, y]$ implique que tous les coefficients $c_{ij}$ sont égaux à
$(nm)!/n!(m!)^n$, comme souhaité.
\end{proof}

\begin{lemma}
\label{lemma-two-ideals}
Soit $A$ un anneau muni de deux idéaux $I, J \subset A$.
Soit $\gamma$ une structure de puissances divisées sur $I$, et soit
$\delta$ une structure de puissances divisées sur $J$.
Alors
\begin{enumerate}
\item $\gamma$ et $\delta$ coïncident sur $IJ$,
\item si $\gamma$ et $\delta$ coïncident sur $I \cap J$, alors elles sont
les restrictions d'une unique structure de puissances divisées $\epsilon$
sur $I + J$.
\end{enumerate}
\end{lemma}

\begin{proof}
Soient $x \in I$ et $y \in J$. Alors
$$
\gamma_n(xy) = y^n\gamma_n(x) = n! \delta_n(y) \gamma_n(x) =
\delta_n(y) x^n = \delta_n(xy).
$$
Par conséquent, $\gamma$ et $\delta$ coïncident sur un ensemble de générateurs
(additifs) de $IJ$. D'après la propriété (4) de la définition
\ref{definition-divided-powers}, elles coïncident donc sur tout $IJ$.

\medskip\noindent
Supposons que $\gamma$ et $\delta$ coïncident sur $I \cap J$.
Soit $z \in I + J$. Écrivons $z = x + y$ avec $x \in I$ et $y \in J$.
On pose alors
$$
\epsilon_n(z) = \sum \gamma_i(x)\delta_{n - i}(y)
$$
pour tout $n \geq 1$.
Pour vérifier que cette définition ne dépend pas du choix, supposons que
$z = x' + y'$ soit une autre écriture avec $x' \in I$ et $y' \in J$. Alors
$w = x - x' = y' - y \in I \cap J$. Par conséquent,
\begin{align*}
\sum\nolimits_{i + j = n} \gamma_i(x)\delta_j(y)
& =
\sum\nolimits_{i + j = n} \gamma_i(x' + w)\delta_j(y) \\
& =
\sum\nolimits_{i' + l + j = n} \gamma_{i'}(x')\gamma_l(w)\delta_j(y) \\
& =
\sum\nolimits_{i' + l + j = n} \gamma_{i'}(x')\delta_l(w)\delta_j(y) \\
& =
\sum\nolimits_{i' + j' = n} \gamma_{i'}(x')\delta_{j'}(y + w) \\
& =
\sum\nolimits_{i' + j' = n} \gamma_{i'}(x')\delta_{j'}(y')
\end{align*}
comme souhaité. Nous avons donc défini des applications
$\epsilon_n : I + J \to I + J$ pour tout $n \geq 1$; il est immédiat que
$\epsilon_n \mid_{I} = \gamma_n$ et
$\epsilon_n \mid_{J} = \delta_n$.
Montrons ensuite les conditions (1) -- (5) de la définition
\ref{definition-divided-powers} pour la famille d'applications $\epsilon_n$.
Les propriétés (1) et (3) sont immédiates. Pour vérifier (4), supposons
que $z = x + y$ et $z' = x' + y'$ avec $x, x' \in I$ et $y, y' \in J$,
et calculons
\begin{align*}
\epsilon_n(z + z') & =
\sum\nolimits_{a + b = n} \gamma_a(x + x')\delta_b(y + y') \\
& =
\sum\nolimits_{i + i' + j + j' = n}
\gamma_i(x) \gamma_{i'}(x')\delta_j(y)\delta_{j'}(y') \\
& =
\sum\nolimits_{k = 0, \ldots, n}
\sum\nolimits_{i+j=k} \gamma_i(x)\delta_j(y)
\sum\nolimits_{i'+j'=n-k} \gamma_{i'}(x')\delta_{j'}(y') \\
& =
\sum\nolimits_{k = 0, \ldots, n}\epsilon_k(z)\epsilon_{n-k}(z')
\end{align*}
ce qui donne le résultat voulu. Il suffit maintenant de vérifier (2) et (5)
sur les éléments de $I$ ou de $J$, voir le lemme
\ref{lemma-check-on-generators}. Cela est immédiat puisque $\gamma$ et
$\delta$ sont des structures de puissances divisées.

\medskip\noindent
Nous avons ainsi démontré l'existence d'une structure de puissances divisées
$\epsilon$ sur $I+J$ dont les restrictions à $I$ et à $J$ sont
$\gamma$ et $\delta$; son unicité est immédiate.
\end{proof}

\begin{lemma}
\label{lemma-nil}
Soit $p$ un nombre premier. Soit $A$ un anneau, soit $I \subset A$ un idéal,
et soit $\gamma$ une structure de puissances divisées sur $I$. Supposons que
$p$ soit nilpotent dans $A/I$. Alors $I$ est localement nilpotent si et
seulement si $p$ est nilpotent dans $A$.
\end{lemma}

\begin{proof}
Si $p^N = 0$ dans $A$, alors, pour $x \in I$, on a
$x^{pN} = (pN)!\gamma_{pN}(x) = 0$, car $(pN)!$ est divisible par $p^N$.
Réciproquement, supposons que $I$ soit localement nilpotent. Comme nous avons
aussi supposé que $p$ est nilpotent dans $A/I$, on a $p^r \in I$ pour un
certain $r$; ainsi $p^r$ est nilpotent, et donc $p$ est nilpotent.
\end{proof}



\section{Anneaux à puissances divisées}

\label{section-divided-power-rings}

\noindent
Les anneaux à puissances divisées forment une catégorie. En voici la définition.

\begin{definition}

\label{definition-divided-power-ring}
Un {\it anneau à puissances divisées} est un triplet $(A, I, \gamma)$, où $A$ est un anneau, $I \subset A$ est un idéal et $\gamma = (\gamma_n)_{n \geq 1}$ est une structure de puissances divisées sur $I$. Un {\it homomorphisme d'anneaux à puissances divisées} $\varphi : (A, I, \gamma) \to (B, J, \delta)$ est un homomorphisme d'anneaux $\varphi : A \to B$ tel que $\varphi(I) \subset J$ et $\delta_n(\varphi(x)) = \varphi(\gamma_n(x))$ pour tout $x \in I$ et tout $n \geq 1$.
\end{definition}

\noindent
Nous dirons parfois « soit $(B, J, \delta)$ une algèbre à puissances divisées sur $(A, I, \gamma)$ » pour indiquer que $(B, J, \delta)$ est un anneau à puissances divisées muni d'un homomorphisme d'anneaux à puissances divisées $(A, I, \gamma) \to (B, J, \delta)$.

\begin{lemma}

\label{lemma-limits}

La catégorie des anneaux à puissances divisées admet toutes les limites, et celles-ci coïncident avec les limites calculées dans la catégorie des anneaux.

\end{lemma}

\begin{proof}
La limite du diagramme vide est l'anneau nul (ce qui est étrange, mais nécessaire ici). Le produit d'une famille d'anneaux à puissances divisées $(A_t, I_t, \gamma_t)$, $t \in T$, est $(\prod A_t, \prod I_t, \gamma)$, où $\gamma_n((x_t)) = (\gamma_{t, n}(x_t))$. L'égalisateur de $\alpha, \beta : (A, I, \gamma) \to (B, J, \delta)$ est simplement $C = \{a \in A \mid \alpha(a) = \beta(a)\}$, muni de l'idéal $C \cap I$ et des puissances divisées induites. Il en résulte que toutes les limites existent ; voir Catégories, lemme \ref{categories-lemma-limits-products-equalizers}.
\end{proof}

\noindent
Le lemme suivant illustre, dans le cas des algèbres à puissances divisées, un phénomène très général de théorie des catégories.

\begin{lemma}

\label{lemma-a-version-of-brown}

Soit $\mathcal{C}$ la catégorie des anneaux à puissances divisées. Soit
$F : \mathcal{C} \to \textit{Sets}$ un foncteur. Supposons que

\begin{enumerate}

\item il existe un cardinal $\kappa$ tel que, pour tout
$f \in F(A, I, \gamma)$, il existe un morphisme
$(A', I', \gamma') \to (A, I, \gamma)$ de $\mathcal{C}$ tel que $f$
soit l'image d'un élément $f' \in F(A', I', \gamma')$ et que $|A'| \leq \kappa$ ;
\item  $F$ commute aux limites.

\end{enumerate}

Alors $F$ est représentable, c'est-à-dire qu'il existe un objet $(B, J, \delta)$
de $\mathcal{C}$ tel que
$$
F(A, I, \gamma) = \Hom_\mathcal{C}((B, J, \delta), (A, I, \gamma))
$$
fonctoriellement en $(A, I, \gamma)$.

\end{lemma}

\begin{proof}
C'est un cas particulier de Catégories, lemme \ref{categories-lemma-a-version-of-brown}.
\end{proof}

\begin{lemma}

\label{lemma-colimits}

La catégorie des anneaux à puissances divisées admet toutes les colimites.

\end{lemma}

\begin{proof}

La colimite du diagramme vide est $\mathbf{Z}$, muni de l'idéal à puissances divisées $(0)$. Considérons maintenant une colimite quelconque. Soit $\mathcal{C}$ une catégorie et soit
$c \mapsto (A_c, I_c, \gamma_c)$ un diagramme. Considérons le foncteur
$$
F(B, J, \delta) = \lim_{c \in \mathcal{C}}
Hom((A_c, I_c, \gamma_c), (B, J, \delta))
$$
Tout élément $f = (f_c)_{c \in C} \in F(B, J, \delta)$ possède la propriété suivante : les images $f_c(A_c)$ engendrent un sous-anneau $B'$ de $B$ dont le cardinal est borné par un cardinal $\kappa$, et les images $f_c(I_c)$ engendrent un sous-idéal à puissances divisées $J'$ de $B'$. On obtient ainsi une factorisation de
$f$ par un élément $f'$ de $F(B')$, suivie de l'inclusion $B' \to B$. En outre,
$F$ commute aux limites. Nous pouvons donc appliquer le lemme \ref{lemma-a-version-of-brown}
pour conclure que $F$ est représentable.

\end{proof}

\begin{remark}

\label{remark-pushout}

Soit
$$
\xymatrix{
(B, J, \delta) \ar[r] & (B'', J'', \delta'') \\
(A, I, \gamma) \ar[r] \ar[u] & (B', J', \delta') \ar[u]
}
$$
un carré cocartésien dans la catégorie des anneaux à puissances divisées (les sommes amalgamées existent d'après le lemme \ref{lemma-colimits}). Alors

\begin{enumerate}

\item  $B''/J'' = B/J \otimes_{A/I} B'/J'$,
\item l'application $\varphi : B \otimes_A B' \to B''$ est surjective, et
\item l'application $J \otimes_A B' + B \otimes_A J' \to J''$ induite par $\varphi$
est surjective.

\end{enumerate}

Pour démontrer (1), considérons les applications $(B'', J'', \delta'') \to (C, (0), \emptyset)$
et appliquons la définition d'un carré cocartésien. Pour démontrer (2), considérons l'image $\tilde B = \varphi(B \otimes_A B')$ et l'idéal
$\tilde J = \varphi(J \otimes_A B' + B \otimes_A J')$. Alors il est clair que $\delta''_n(\tilde J) \subset \tilde J$ pour tous $n \geq 1$
par compatibilité de $\delta''$ avec $\delta$ et $\delta'$. La propriété universelle du carré cocartésien impose alors $\tilde B = B''$ et $\tilde J = J''$.

\end{remark}

\begin{remark}

\label{remark-forgetful}

Le foncteur d'oubli $(A, I, \gamma) \mapsto A$ ne commute pas aux colimites. Par exemple, soit
$$
\xymatrix{
(B, J, \delta) \ar[r] & (B'', J'', \delta'') \\
(A, I, \gamma) \ar[r] \ar[u] & (B', J', \delta') \ar[u]
}
$$
un carré cocartésien dans la catégorie des anneaux à puissances divisées, comme dans la remarque \ref{remark-pushout}. En général, l'application
$B \otimes_A B' \to B''$ n'est pas un isomorphisme. Un exemple explicite est fourni par
$(A, I, \gamma) = (\mathbf{Z}, (0), \emptyset)$,
$(B, J, \delta) = (\mathbf{Z}/4\mathbf{Z}, 2\mathbf{Z}/4\mathbf{Z}, \delta)$,
$(B', J', \delta') =
(\mathbf{Z}/4\mathbf{Z}, 2\mathbf{Z}/4\mathbf{Z}, \delta')$
où $\delta_2(2) = 2$ et $\delta'_2(2) = 0$. Plus précisément, en appliquant le lemme \ref{lemma-need-only-gamma-p},
on prend pour $\delta$, respectivement\ $\delta'$, l'unique structure de puissances divisées sur $J$, respectivement\ $J'$, telle que
$\delta_2 : J \to J$, respectivement\ $\delta'_2 : J' \to J'$, soit l'application $0 \mapsto 0, 2 \mapsto 2$, respectivement\ $0 \mapsto 0, 2 \mapsto 0$. Alors $(B'', J'', \delta'') = (\mathbf{F}_2, (0), \emptyset)$,
ce qui ne coïncide pas avec le produit tensoriel.

\end{remark}

\section{Prolongement des puissances divisées}

\label{section-extend}

\noindent
Voici la définition.

\begin{definition}

\label{definition-extends}
Étant donnés un anneau à puissances divisées $(A, I, \gamma)$ et un homomorphisme d'anneaux $A \to B$, nous disons que $\gamma$ {\it se prolonge} à $B$ s'il existe une structure de puissances divisées $\bar \gamma$ sur $IB$ telle que $(A, I, \gamma) \to (B, IB, \bar\gamma)$ soit un homomorphisme d'anneaux à puissances divisées.
\end{definition}

\begin{lemma}

\label{lemma-gamma-extends}
Soit $(A, I, \gamma)$ un anneau à puissances divisées. Soit $A \to B$ un homomorphisme d'anneaux. Si $\gamma$ se prolonge à $B$, ce prolongement est unique. Supposons que l'une au moins des conditions suivantes soit vérifiée :
\begin{enumerate}
\item $IB = 0$,
\item $I$ est principal, ou
\item $A \to B$ est plat.
\end{enumerate}
Alors $\gamma$ se prolonge à $B$.
\end{lemma}

\begin{proof}
Tout élément de $IB$ s'écrit comme une somme finie $\sum\nolimits_{i=1}^t b_ix_i$, avec $b_i \in B$ et $x_i \in I$. Si $\gamma$ admet un prolongement $\bar\gamma$ à $IB$, alors $\bar\gamma_n(x_i) = \gamma_n(x_i)$. Les conditions (3) et (4) de la définition \ref{definition-divided-powers} impliquent donc que
$$
\bar\gamma_n(\sum\nolimits_{i=1}^t b_ix_i) =
\sum\nolimits_{n_1 + \ldots + n_t = n}
\prod\nolimits_{i = 1}^t b_i^{n_i}\gamma_{n_i}(x_i)
$$
On voit ainsi que $\bar\gamma$, s'il existe, est unique.

\medskip\noindent

Si $IB = 0$, il suffit de poser $\bar\gamma_n(0) = 0$. Si $I = (x)$,
nous définissons $\bar\gamma_n(bx) = b^n\gamma_n(x)$. Cette définition ne dépend pas du coefficient choisi : si $b'x = bx$, c'est-à-dire si $(b - b')x = 0$, alors

\begin{align*}
b^n\gamma_n(x) - (b')^n\gamma_n(x)
& =
(b^n - (b')^n)\gamma_n(x) \\
& =
(b^{n - 1} + \ldots + (b')^{n - 1})(b - b')\gamma_n(x) = 0
\end{align*}

car $\gamma_n(x)$ est divisible par $x$ (puisque
$\gamma_n(I) \subset I$) et est donc annulé par $b - b'$. Vérifions maintenant les conditions (1) -- (5) de la définition \ref{definition-divided-powers}. Les conditions (1), (2), (3) et (5) découlent immédiatement de la construction. Pour (4), soient $y, z \in IB$, disons $y = bx$ et $z = cx$. Alors
$y + z = (b + c)x$. Par conséquent,

\begin{align*}
\bar\gamma_n(y + z)
& =
(b + c)^n\gamma_n(x) \\
& =
\sum \frac{n!}{i!(n - i)!}b^ic^{n -i}\gamma_n(x) \\
& =
\sum b^ic^{n - i}\gamma_i(x)\gamma_{n - i}(x) \\
& =
\sum \bar\gamma_i(y)\bar\gamma_{n -i}(z)
\end{align*}

comme voulu.

\medskip\noindent

Supposons enfin que $A \to B$ soit plat. Soient $b_1, \ldots, b_r \in B$ et
$x_1, \ldots, x_r \in I$. Alors
$$
\bar\gamma_n(\sum b_ix_i) =
\sum b_1^{e_1} \ldots b_r^{e_r} \gamma_{e_1}(x_1) \ldots \gamma_{e_r}(x_r)
$$
où la somme porte sur les $e_1 + \ldots + e_r = n$,
si $\bar\gamma_n$ existe. Supposons ensuite donnés $c_1, \ldots, c_s \in B$
et $a_{ij} \in A$ tels que $b_i = \sum a_{ij}c_j$. En posant $y_j = \sum a_{ij}x_i$, nous affirmons que
$$
\sum b_1^{e_1} \ldots b_r^{e_r} \gamma_{e_1}(x_1) \ldots \gamma_{e_r}(x_r) =
\sum c_1^{d_1} \ldots c_s^{d_s} \gamma_{d_1}(y_1) \ldots \gamma_{d_s}(y_s)
$$
dans $B$, où, au membre de droite, la somme porte sur les $d_1 + \ldots + d_s = n$. En effet, les axiomes d'une structure de puissances divisées permettent de développer les deux membres en sommes, à coefficients dans $\mathbf{Z}[a_{ij}]$, de termes de la forme
$c_1^{d_1} \ldots c_s^{d_s}\gamma_{e_1}(x_1) \ldots \gamma_{e_r}(x_r)$. Pour vérifier l'égalité des coefficients, il suffit de remarquer que le résultat est vrai dans $\mathbf{Q}[x_1, \ldots, x_r, c_1, \ldots, c_s, a_{ij}]$, muni de
l'unique structure de puissances divisées $\gamma$ sur $(x_1, \ldots, x_r)$. D'après le théorème de Lazard (Algèbre, théorème \ref{algebra-theorem-lazard}), $B$ est une colimite filtrante de $A$-modules libres de type fini. En particulier, si $z \in IB$ s'écrit à la fois $z = \sum x_ib_i$ et
$z = \sum x'_{i'}b'_{i'}$, alors nous pouvons trouver $c_1, \ldots, c_s \in B$
et $a_{ij}, a'_{i'j} \in A$ tels que $b_i = \sum a_{ij}c_j$
et $b'_{i'} = \sum a'_{i'j}c_j$, et tels que
$y_j = \sum x_ia_{ij} = \sum x'_{i'}a'_{i'j}$\footnote{On peut également le démontrer sans recourir au théorème \ref{algebra-theorem-lazard} d'Algèbre. En effet, si
$z = \sum x_ib_i$ et $z = \sum x'_{i'}b'_{i'}$, alors
$\sum x_ib_i - \sum x'_{i'}b'_{i'} = 0$ est une relation dans le
$A$-module $B$. Le lemme \ref{algebra-lemma-flat-eq} d'Algèbre
(appliqué aux $x_i$ et $x'_{i'}$ à la place des $f_i$, et aux $b_i$ et $b'_{i'}$ à la place des $x_i$) fournit alors les éléments $c_1, \ldots, c_s \in B$
et $a_{ij}, a'_{i'j} \in A$ voulus.}. La procédure ci-dessus définit donc une application $\bar\gamma_n$
sur $IB$, indépendante de l'écriture choisie. Par construction, $\bar\gamma$ satisfait les conditions (1), (3) et (4). De plus, pour $x \in I$, on a $\bar\gamma_n(x) = \gamma_n(x)$. Le lemme \ref{lemma-check-on-generators} montre alors que $\bar\gamma$
est une structure de puissances divisées sur $IB$.

\end{proof}

\begin{lemma}

\label{lemma-kernel}

Soit $(A, I, \gamma)$ un anneau à puissances divisées.

\begin{enumerate}

\item Si $\varphi : (A, I, \gamma) \to (B, J, \delta)$ est un homomorphisme d'anneaux à puissances divisées, alors $\Ker(\varphi) \cap I$
est préservé par $\gamma_n$ pour tous $n \geq 1$.
\item Soit $\mathfrak a \subset A$ un idéal et posons
$I' = I \cap \mathfrak a$. Les conditions suivantes sont équivalentes :

\begin{enumerate}

\item  $I'$ est préservé par $\gamma_n$ pour tous $n > 0$,
\item  $\gamma$ se prolonge à $A/\mathfrak a$,
\item il existe une famille de générateurs $x_i$ de l'idéal $I'$ telle que $\gamma_n(x_i) \in I'$ pour tout $n > 0$.

\end{enumerate}

\end{enumerate}

\end{lemma}

\begin{proof}

Démonstration de (1). C'est immédiat. Supposons (2)(a). Posons
$\bar\gamma_n(x \bmod I') = \gamma_n(x) \bmod I'$ pour $x \in I$. Cette définition est indépendante du représentant, car on a $\gamma_n(x + y) = \gamma_n(x) \bmod I'$
pour $y \in I'$
d'après la définition \ref{definition-divided-powers}, (4), et l'hypothèse $\gamma_j(y) \in I'$. Il est clair que $\bar\gamma$ est une structure de puissances divisées puisque $\gamma$ en est une. Ainsi, (2)(b) est vérifiée. Réciproquement, (2)(b) implique (2)(a) d'après (1). Il est clair que (2)(a) implique (2)(c). Supposons (2)(c). Notons que $\gamma_n(x) = a^n\gamma_n(x_i) \in I'$ lorsque $x = ax_i$. Ainsi, $\gamma_n(x) \in I'$ pour une famille de générateurs de $I'$ comme groupe abélien. Par récurrence sur la longueur d'une expression en ces générateurs, il suffit de montrer que $\forall n : \gamma_n(x + y) \in I'$
dès que $\forall n : \gamma_n(x), \gamma_n(y) \in I'$. Cela résulte immédiatement du quatrième axiome d'une structure de puissances divisées.

\end{proof}

\begin{lemma}

\label{lemma-sub-dp-ideal}
Soit $(A, I, \gamma)$ un anneau à puissances divisées et soit $E \subset I$ une partie. Le plus petit idéal $J \subset I$ stable par $\gamma$ et contenant tous les $f \in E$ est l'idéal $J$ engendré par les $\gamma_n(f)$, pour $n \geq 1$ et $f \in E$.
\end{lemma}

\begin{proof}
Cela résulte immédiatement du lemme \ref{lemma-kernel}.
\end{proof}

\begin{lemma}

\label{lemma-extend-to-completion}

Soit $(A, I, \gamma)$ un anneau à puissances divisées. Soit $p$ un nombre premier. Si $p$ est nilpotent dans $A/I$, alors

\begin{enumerate}

\item l'application canonique du complété $p$-adique $A^\wedge = \lim_e A/p^eA$ vers $A/I$ est surjective ;
\item le noyau de cette application est le complété $p$-adique $I^\wedge$ de $I$ ;
\item chaque application $\gamma_n$ est continue pour la topologie $p$-adique et se prolonge en une application $\gamma_n^\wedge : I^\wedge \to I^\wedge$, les applications ainsi obtenues définissant une structure de puissances divisées sur $I^\wedge$.

\end{enumerate}

Si, de plus, $A$ est une $\mathbf{Z}_{(p)}$-algèbre, alors

\begin{enumerate}

\item[(4)] pour $e$ assez grand, l'idéal $p^eA \subset I$ est stable par la structure de puissances divisées $\gamma$, et
$$
(A^\wedge, I^\wedge, \gamma^\wedge) = \lim_e (A/p^eA, I/p^eA, \bar\gamma)
$$
dans la catégorie des anneaux à puissances divisées.

\end{enumerate}

\end{lemma}

\begin{proof}

Soit $t \geq 1$ un entier tel que $p^tA/I = 0$, c'est-à-dire $p^tA \subset I$. L'application $A^\wedge \to A/I$ est la composée $A^\wedge \to A/p^tA \to A/I$,
qui est surjective (par exemple d'après Algèbre, lemme \ref{algebra-lemma-completion-generalities}). Comme $p^eI \subset p^eA \cap I \subset p^{e - t}I$ pour $e \geq t$, le noyau de la composée $A^\wedge \to A/I$ est le complété $p$-adique de $I$. L'application $\gamma_n$ est continue, car
$$
\gamma_n(x + p^ey) =
\sum\nolimits_{i + j = n} p^{je}\gamma_i(x)\gamma_j(y) =
\gamma_n(x) \bmod p^eI
$$
par les axiomes d'une structure de puissances divisées. Il est clair que les applications $\gamma_n^\wedge$ satisfont aux mêmes axiomes que les applications $\gamma_n$.
Enfin, pour démontrer la dernière assertion, prenons $e > t$. Alors $p^eA \subset I$ et $\gamma_1(p^eA) \subset p^eA$ ; pour $n > 1$, on a
$$
\gamma_n(p^ea) = p^n \gamma_n(p^{e - 1}a) = \frac{p^n}{n!} p^{n(e - 1)}a^n
\in p^e A
$$
car $p^n/n! \in \mathbf{Z}_{(p)}$ et, puisque $n \geq 2$ et $e \geq 2$,
$n(e - 1) \geq e$. Cela prouve que $\gamma$ se prolonge à $A/p^eA$ ; voir le lemme \ref{lemma-kernel}. L'assertion relative aux limites découle immédiatement de la construction des limites donnée dans la démonstration du lemme \ref{lemma-limits}.

\end{proof}
\section{Algèbres polynomiales à puissances divisées}
\label{section-divided-power-polynomial-ring}

\noindent
Un exemple très utile est celui de l'{\it algèbre polynomiale à puissances
divisées}. Soit $A$ un anneau. Soit $t \geq 1$. On note
$A\langle x_1, \ldots, x_t \rangle$ la $A$-algèbre suivante :
en tant que $A$-module, on pose
$$
A\langle x_1, \ldots, x_t \rangle =
\bigoplus\nolimits_{n_1, \ldots, n_t \geq 0} A x_1^{[n_1]} \ldots x_t^{[n_t]}
$$
et la multiplication est donnée par
$$
x_i^{[n]}x_i^{[m]} = \frac{(n + m)!}{n!m!}x_i^{[n + m]}.
$$
On pose aussi $x_i = x_i^{[1]}$. Remarquons que
$1 = x_1^{[0]} \ldots x_t^{[0]}$. Une construction analogue donne
l'algèbre polynomiale à puissances divisées en une infinité de variables.
Il existe un morphisme canonique de $A$-algèbres
$A\langle x_1, \ldots, x_t \rangle \to A$ qui envoie $x_i^{[n]}$ sur zéro
pour $n > 0$. Le noyau de ce morphisme est noté
$A\langle x_1, \ldots, x_t \rangle_{+}$.

\begin{lemma}
\label{lemma-divided-power-polynomial-algebra}
Soit $(A, I, \gamma)$ un anneau à puissances divisées.
Il existe une unique structure de puissances divisées $\delta$ sur
$$
J = IA\langle x_1, \ldots, x_t \rangle + A\langle x_1, \ldots, x_t \rangle_{+}
$$
telle que
\begin{enumerate}
\item $\delta_n(x_i) = x_i^{[n]}$, et
\item $(A, I, \gamma) \to (A\langle x_1, \ldots, x_t \rangle, J, \delta)$
soit un homomorphisme d'anneaux à puissances divisées.
\end{enumerate}
De plus, $(A\langle x_1, \ldots, x_t \rangle, J, \delta)$ possède la
propriété universelle suivante : se donner un homomorphisme d'anneaux à
puissances divisées
$\varphi : (A\langle x_1, \ldots, x_t \rangle, J, \delta) \to
(C, K, \epsilon)$ revient à se donner un homomorphisme d'anneaux à
puissances divisées $A \to C$ et des éléments
$k_1, \ldots, k_t \in K$.
\end{lemma}

\begin{proof}
Nous démontrons le lemme dans le cas d'une algèbre polynomiale à puissances
divisées en une variable. Le cas général se traite exactement de la même
manière ; on peut aussi remarquer que $A\langle x_1, \ldots, x_t\rangle$
est isomorphe à l'anneau obtenu en adjoignant successivement les variables à
puissances divisées $x_1, \ldots, x_t$.

\medskip\noindent
Soit $A\langle x \rangle_{+}$ l'idéal engendré par
$x, x^{[2]}, x^{[3]}, \ldots$. Remarquons que
$J = IA\langle x \rangle + A\langle x \rangle_{+}$ et que
$$
IA\langle x \rangle \cap A\langle x \rangle_{+} =
IA\langle x \rangle \cdot A\langle x \rangle_{+}
$$
Par conséquent, d'après le lemme \ref{lemma-two-ideals}, il suffit de montrer
qu'il existe des structures de puissances divisées sur les idéaux
$IA\langle x \rangle$ et $A\langle x \rangle_{+}$. L'existence de la
première résulte du lemme \ref{lemma-gamma-extends}, puisque
$A \to A\langle x \rangle$ est plat. Pour la seconde, si $A$ est sans
torsion, le lemme \ref{lemma-silly} (4) montre que $\delta$ existe. En effet,
en choisissant les éléments $x^{[m]}$ comme générateurs, on constate que
$(x^{[m]})^n = \frac{(nm)!}{(m!)^n} x^{[nm]}$
et que $n!$ divise l'entier $\frac{(nm)!}{(m!)^n}$.
Dans le cas général, écrivons $A = R/\mathfrak a$, où $R$ est un anneau
sans torsion (par exemple un anneau de polynômes sur $\mathbf{Z}$). Le noyau
de $R\langle x \rangle \to A\langle x \rangle$ est
$\bigoplus \mathfrak a x^{[m]}$. En appliquant le critère (2)(c) du
lemme \ref{lemma-kernel}, on voit que la structure de puissances divisées
sur $R\langle x \rangle_{+}$ se transporte à $A\langle x \rangle$, comme
souhaité.

\medskip\noindent
Démontrons la propriété universelle. Étant donnés un homomorphisme
$\varphi : A \to C$ d'anneaux à puissances divisées et des éléments
$k_1, \ldots, k_t \in K$, considérons
$$
A\langle x_1, \ldots, x_t \rangle \to C,\quad
x_1^{[n_1]} \ldots x_t^{[n_t]} \longmapsto
\epsilon_{n_1}(k_1) \ldots \epsilon_{n_t}(k_t)
$$
où l'on utilise $\varphi$ sur les coefficients. Il reste seulement à
vérifier qu'il s'agit d'un homomorphisme de $A$-algèbres (détails omis).
La construction inverse est claire.
\end{proof}

\begin{remark}
\label{remark-divided-power-polynomial-algebra}
Soit $(A, I, \gamma)$ un anneau à puissances divisées. Il existe une variante
du lemme \ref{lemma-divided-power-polynomial-algebra} pour une infinité de
variables. Remarquons d'abord que, si $s < t$, il existe un morphisme
canonique
$$
A\langle x_1, \ldots, x_s \rangle \to A\langle x_1, \ldots, x_t\rangle
$$
Ainsi, si $W$ est un ensemble quelconque, on pose
$$
A\langle x_w: w \in W\rangle =
\colim_{E \subset W} A\langle x_e:e \in E\rangle
$$
(la colimite étant prise sur les sous-ensembles finis $E$ de $W$), avec les
morphismes de transition ci-dessus. La définition d'une colimite montre que
la propriété universelle de $A\langle x_w: w \in W\rangle$ est entièrement
analogue à celle énoncée dans le
lemme \ref{lemma-divided-power-polynomial-algebra}.
\end{remark}

\noindent
Le lemme suivant se trouve dans \cite{BO}.

\begin{lemma}
\label{lemma-need-only-gamma-p}
Soit $p$ un nombre premier. Soit $A$ un anneau dans lequel tout entier $n$
non divisible par $p$ est inversible, autrement dit $A$ est une
$\mathbf{Z}_{(p)}$-algèbre. Soit $I \subset A$ un idéal. Deux structures de
puissances divisées $\gamma, \gamma'$ sur $I$ sont égales si et seulement si
$\gamma_p = \gamma'_p$. De plus, étant donnée une application
$\delta : I \to I$ telle que
\begin{enumerate}
\item $p!\delta(x) = x^p$ pour tout $x \in I$\footnote{Cette condition
résulte des deux autres ; voir la remarque \ref{remark-ryo-suzuki}.},
\item $\delta(ax) = a^p\delta(x)$ pour tous $a \in A$, $x \in I$, et
\item
$\delta(x + y) =
\delta(x) +
\sum\nolimits_{i + j = p, i,j \geq 1} \frac{1}{i!j!} x^i y^j +
\delta(y)$ pour tous $x, y \in I$,
\end{enumerate}
il existe une unique structure de puissances divisées $\gamma$ sur $I$ telle
que $\gamma_p = \delta$.
\end{lemma}

\begin{proof}
Si $n$ n'est pas divisible par $p$, alors
$\gamma_n(x) = c x \gamma_{n - 1}(x)$, où $c$ est une unité de
$\mathbf{Z}_{(p)}$. De plus,
$$
\gamma_{pm}(x) = c \gamma_m(\gamma_p(x))
$$
où $c$ est une unité de $\mathbf{Z}_{(p)}$. La première assertion est donc
claire. Pour la seconde, on peut remonter ces égalités afin de définir tous
les $\gamma_n$. Plus précisément, si
$n = a_0 + a_1p + \ldots + a_e p^e$ avec $a_i \in \{0, \ldots, p - 1\}$,
on pose
$$
\gamma_n(x) = c_n x^{a_0} \delta(x)^{a_1} \ldots \delta^e(x)^{a_e}
$$
où $c_n \in \mathbf{Z}_{(p)}$ est défini par
$$
c_n =
{(p!)^{a_1 + a_2(1 + p) + \ldots + a_e(1 + \ldots + p^{e - 1})}}/{n!}.
$$
Il faut maintenant vérifier les axiomes (1) -- (5) d'une structure de
puissances divisées ; voir la définition
\ref{definition-divided-powers}. Les axiomes (1) et (3) sont immédiats.
Les axiomes (2) et (5) se vérifient par un calcul direct que nous omettons.
Soient $x, y \in I$. Nous affirmons qu'il existe un homomorphisme d'anneaux
$$
\varphi : \mathbf{Z}_{(p)}\langle u, v \rangle \longrightarrow A
$$
qui envoie $u^{[n]}$ sur $\gamma_n(x)$ et $v^{[n]}$ sur $\gamma_n(y)$.
D'après la construction de $\mathbf{Z}_{(p)}\langle u, v \rangle$, il faut
pour cela vérifier que
$$
\gamma_n(x)\gamma_m(x) = \frac{(n + m)!}{n!m!} \gamma_{n + m}(x)
$$
dans $A$, et de même pour $y$. C'est vrai, puisque l'axiome (2) est vérifié
par $\gamma$. Notons $\epsilon$ la structure de puissances divisées sur
l'idéal $\mathbf{Z}_{(p)}\langle u, v\rangle_{+}$ de
$\mathbf{Z}_{(p)}\langle u, v\rangle$. Nous affirmons ensuite que
$\varphi(\epsilon_n(f)) = \gamma_n(\varphi(f))$
pour $f \in \mathbf{Z}_{(p)}\langle u, v\rangle_{+}$ et tout $n$.
C'est clair pour $n = 0, 1, \ldots, p - 1$. Pour $n = p$, il suffit de le
vérifier sur un système de générateurs de l'idéal
$\mathbf{Z}_{(p)}\langle u, v\rangle_{+}$, car $\epsilon_p$ et
$\gamma_p = \delta$ vérifient toutes deux les propriétés (1) et (3) du
lemme. Il suffit donc d'établir
$\gamma_p(\gamma_n(x)) = \frac{(pn)!}{p!(n!)^p}\gamma_{pn}(x)$, et de même
pour $y$ ; cela résulte de l'axiome (5) pour $\gamma$.
Maintenant, si $n = a_0 + a_1p + \ldots + a_e p^e$ est un entier quelconque
écrit sous la forme de son développement $p$-adique ci-dessus, alors
$$
\epsilon_n(f) =
c_n f^{a_0} \gamma_p(f)^{a_1} \ldots \gamma_p^e(f)^{a_e}
$$
car $\epsilon$ est une structure de puissances divisées. On en déduit que
$\varphi(\epsilon_n(f)) = \gamma_n(\varphi(f))$ pour tout $n$. En appliquant
ceci à $f = u + v$, on voit que l'axiome (4) pour $\gamma$ découle du fait
que $\epsilon$ est une structure de puissances divisées.
\end{proof}

\begin{remark}
\label{remark-ryo-suzuki}
Dans la situation $A \supset I$ du lemme
\ref{lemma-need-only-gamma-p}, soit $\delta : I \to I$ une application
vérifiant (2) et (3) du lemme. Alors (1) est également vérifiée.
Montrons d'abord que
$\delta(nx) = n\delta(x) + \frac{n^p - n}{p!}x^p$ pour $n \geq 1$, par
récurrence sur $n$. Le cas $n = 1$ est clair. Supposons le résultat établi
pour un certain $n$.
Alors
$\delta((n + 1)x) = \delta(nx) + \delta(x) +
(\sum_{i = 1}^{p - 1} \frac{n^i}{i!(p - i)!})x^p
= (n + 1)\delta(x) + \left( \frac{n^p - n}{p!} + \frac{(n + 1)^p}{p!}
- \frac{n^p + 1}{p!} \right)x^p =
(n + 1)\delta(x) + \frac{(n + 1)^p - (n + 1)}{p!}x^p$,
ce qui établit le résultat pour $n + 1$. D'autre part,
$\delta(nx) = n^p\delta(x)$, donc
$(n^p - n)\delta(x) = \frac{n^p - n}{p!}x^p$.
Or, pour tout $p$, il existe un entier $n$ tel que
$p^2 \not \mid n^p - n$. En effet,
$(\mathbf{Z}/p^2\mathbf{Z})^\times \cong \mathbf{Z}/p(p-1)\mathbf{Z}$.
Il existe donc $n$ tel que $n^{p-1} \not \equiv 1 \mod p^2$.
On obtient ainsi $p!\delta(x)=x^p$.
\end{remark}









\section{Résolutions de Tate}
\label{section-tate}

\noindent
Dans cette section, nous présentons brièvement les résolutions construites
dans \cite{Tate-homology} et \cite{AH}, qui combinent les structures de
puissances divisées et les algèbres différentielles graduées. Nous utiliserons
la {\it notation homologique} pour les algèbres différentielles graduées.
Celles-ci seront concentrées en degrés homologiques positifs ou nuls. Ainsi,
nos algèbres différentielles graduées $(A, \text{d})$ seront des complexes de
chaînes
$$
\ldots \to A_2 \to A_1 \to A_0 \to 0 \to \ldots
$$
munis d'une multiplication.

\medskip\noindent
Soit $R$ un anneau (commutatif, comme toujours). Dans cette section, nous
considérerons souvent des $R$-algèbres graduées
$A = \bigoplus_{d \geq 0} A_d$ dont les composantes de degré négatif sont
nulles. On pose $A_+ = \bigoplus_{d > 0} A_d$. On notera
$A_{even} = \bigoplus_{d \geq 0} A_{2d}$ et
$A_{odd} = \bigoplus_{d \geq 0} A_{2d + 1}$. Rappelons que $A$ est
commutative au sens gradué si
$x y = (-1)^{\deg(x)\deg(y)} y x$ pour tous éléments homogènes $x, y$.
Rappelons que $A$ est strictement commutative au sens gradué si, de plus,
$x^2 = 0$ pour tout élément homogène $x$ de degré impair. Enfin, pour
comprendre la définition suivante, gardons à l'esprit que
$\gamma_n(x) = x^n/n!$ lorsque $A$ est une $\mathbf{Q}$-algèbre.

\begin{definition}
\label{definition-divided-powers-graded}
Soit $R$ un anneau. Soit $A = \bigoplus_{d \geq 0} A_d$ une $R$-algèbre
graduée strictement commutative au sens gradué. Une famille d'applications
$\gamma_n : A_{even, +} \to A_{even, +}$, définies pour tout $n > 0$, est
appelée une {\it structure de puissances divisées} sur $A$ si les conditions
suivantes sont satisfaites :
\begin{enumerate}
\item $\gamma_n(x) \in A_{2nd}$ si $x \in A_{2d}$,
\item $\gamma_1(x) = x$ pour tout $x$ ; on pose aussi $\gamma_0(x) = 1$,
\item $\gamma_n(x)\gamma_m(x) = \frac{(n + m)!}{n! m!} \gamma_{n + m}(x)$,
\item $\gamma_n(xy) = x^n \gamma_n(y)$ pour tous $x \in A_{even}$ et
$y \in A_{even, +}$,
\item $\gamma_n(xy) = 0$ si $x, y \in A_{odd}$ sont homogènes et $n > 1$,
\item si $x, y \in A_{even, +}$, alors
$\gamma_n(x + y) = \sum_{i = 0, \ldots, n} \gamma_i(x)\gamma_{n - i}(y)$,
\item $\gamma_n(\gamma_m(x)) =
\frac{(nm)!}{n! (m!)^n} \gamma_{nm}(x)$ pour $x \in A_{even, +}$.
\end{enumerate}
\end{definition}

\noindent
Remarquons que les conditions (2), (3), (4), (6) et (7) impliquent que
$\gamma$ est une structure de puissances divisées « usuelle » sur l'idéal
$A_{even, +}$ de l'anneau commutatif $A_{even}$ ; voir les
sections \ref{section-divided-powers},
\ref{section-divided-power-rings},
\ref{section-extend} et
\ref{section-divided-power-polynomial-ring}.
En particulier, $n! \gamma_n(x) = x^n$ pour tout $x \in A_{even, +}$.
La condition (1) exprime la compatibilité de $\gamma$ avec la graduation, et
la condition (5) dit que $\gamma_n$, pour $n > 1$, s'annule sur les produits
d'éléments homogènes de degré impair. Il peut toutefois arriver que
$$
\gamma_2(z_1 z_2 + z_3 z_4) = z_1z_2z_3z_4
$$
soit non nul lorsque $z_1, z_2, z_3, z_4$ sont des éléments homogènes de
degré impair.

\begin{example}[Adjonction d'une variable impaire]
\label{example-adjoining-odd}
Soit $R$ un anneau. Soit $(A, \gamma)$ une $R$-algèbre graduée strictement
commutative au sens gradué, munie d'une structure de puissances divisées
comme dans la définition ci-dessus. Soit $d > 0$ un entier impair. Dans ce
cas, on peut adjoindre une variable $T$ de degré $d$ à $A$. Plus précisément,
on pose
$$
A\langle T \rangle = A \oplus AT
$$
avec la graduation donnée par
$A\langle T \rangle_m = A_m \oplus A_{m - d}T$. Il existe une unique
structure de puissances divisées sur $A\langle T \rangle$ compatible avec la
structure donnée sur $A$. Elle est définie par
$$
\gamma_n(x + yT) = \gamma_n(x) + \gamma_{n - 1}(x)yT
$$
pour $x \in A_{even, +}$ et $y \in A_{odd}$.
\end{example}

\begin{example}[Adjonction d'une variable paire]
\label{example-adjoining-even}
Soit $R$ un anneau. Soit $(A, \gamma)$ une $R$-algèbre graduée strictement
commutative au sens gradué, munie d'une structure de puissances divisées
comme dans la définition ci-dessus. Soit $d > 0$ un entier pair. Dans ce cas,
on peut adjoindre une variable $T$ de degré $d$ à $A$. Plus précisément, on
pose
$$
A\langle T \rangle = A \oplus AT \oplus AT^{(2)} \oplus AT^{(3)} \oplus \ldots
$$
avec la multiplication donnée par
$$
T^{(n)} T^{(m)} = \frac{(n + m)!}{n!m!} T^{(n + m)}
$$
et la graduation donnée par
$$
A\langle T \rangle_m =
A_m \oplus A_{m - d}T \oplus A_{m - 2d}T^{(2)} \oplus \ldots
$$
Il existe une unique structure de puissances divisées sur
$A\langle T \rangle$, compatible avec celle donnée sur $A$, telle que
$\gamma_n(T^{(i)}) = T^{(ni)}$. Pour la définir, on pose d'abord
$$
\gamma_n\left(\sum\nolimits_{i > 0} x_i T^{(i)}\right) =
\sum \prod\nolimits_{n = \sum e_i} x_i^{e_i} T^{(ie_i)}
$$
si $x_i$ appartient à $A_{even}$. Si $x_0 \in A_{even, +}$, on pose ensuite
$$
\gamma_n\left(\sum\nolimits_{i \geq 0} x_i T^{(i)}\right) =
\sum\nolimits_{a + b = n}
\gamma_a(x_0)\gamma_b\left(\sum\nolimits_{i > 0} x_iT^{(i)}\right)
$$
où $\gamma_b$ est défini comme ci-dessus.
\end{example}

\begin{remark}
\label{remark-adjoining-set-of-variables}
On peut aussi adjoindre un ensemble, éventuellement infini, de générateurs
extérieurs ou à puissances divisées dans un degré donné $d > 0$, plutôt qu'un
seul comme dans les exemples \ref{example-adjoining-odd} et
\ref{example-adjoining-even}. Plus précisément, suivant la remarque
\ref{remark-divided-power-polynomial-algebra}, pour $(A,\gamma)$ comme
ci-dessus et un ensemble $J$, soit $A\langle
T_j:j\in J\rangle$ la colimite filtrante des algèbres
$A\langle T_j:j\in S\rangle$, où $S$ parcourt les
sous-ensembles finis de $J$. Il est immédiat que cette algèbre possède une
unique structure de puissances divisées compatible avec la structure donnée
sur $A$ et avec celle de chaque générateur $T_j$.
\end{remark}

\noindent
Nous relions maintenant la notion de structure de puissances divisées à
celle de différentielle. Pour comprendre la définition, remarquons que
$\text{d}(x^n/n!) = \text{d}(x) x^{n - 1}/(n - 1)!$ si $A$ est une
$\mathbf{Q}$-algèbre et si $x \in A_{even, +}$.

\begin{definition}
\label{definition-divided-powers-dga}
Soit $R$ un anneau. Soit $A = \bigoplus_{d \geq 0} A_d$ une $R$-algèbre
différentielle graduée strictement commutative au sens gradué. Une structure
de puissances divisées $\gamma$ sur $A$ est {\it compatible avec la structure
différentielle graduée} si
$\text{d}(\gamma_n(x)) = \text{d}(x) \gamma_{n - 1}(x)$ pour tout
$x \in A_{even, +}$.
\end{definition}

\noindent
Attention : soit $(A, \text{d}, \gamma)$ comme dans la définition
\ref{definition-divided-powers-dga}. Il se peut que $\gamma_n(x)$ ne soit pas
un bord, même lorsque $x$ est un bord. Ainsi, en général, $\gamma$
n'induit pas de structure de puissances divisées sur l'algèbre d'homologie
$H(A)$. Dans certains articles, les auteurs imposent une condition de
compatibilité supplémentaire afin que ce soit le cas ; nous choisissons de
ne pas le faire.

\begin{lemma}
\label{lemma-dpdga-good}
Soient $(A, \text{d}, \gamma)$ et $(B, \text{d}, \gamma)$ comme dans la
définition \ref{definition-divided-powers-dga}. Soit $f : A \to B$ un
morphisme d'algèbres différentielles graduées compatible avec les structures
de puissances divisées. Supposons que
\begin{enumerate}
\item $H_k(A) = 0$ pour $k > 0$, et
\item $f$ soit surjectif.
\end{enumerate}
Alors $\gamma$ induit une structure de puissances divisées sur la
$R$-algèbre graduée $H(B)$.
\end{lemma}

\begin{proof}
Supposons que $x$ et $x'$ soient homogènes de même degré $2d$ et représentent
la même classe d'homologie dans $H(B)$. Écrivons
$x' - x = \text{d}(w)$. Choisissons un relèvement $y \in A_{2d}$ de $x$ et
un relèvement $z \in A_{2d + 1}$ de $w$. Alors
$y' = y + \text{d}(z)$ est un relèvement de $x'$. Par conséquent,
$$
\gamma_n(y') = \sum \gamma_i(y) \gamma_{n - i}(\text{d}(z))
= \gamma_n(y) +
\sum\nolimits_{i < n} \gamma_i(y) \gamma_{n - i}(\text{d}(z))
$$
Comme $A$ est acyclique en degrés positifs et que
$\text{d}(\gamma_j(\text{d}(z))) = 0$ pour tout $j$, on peut écrire cette
égalité sous la forme
$$
\gamma_n(y') = \gamma_n(y) +
\sum\nolimits_{i < n} \gamma_i(y) \text{d}(z_i)
$$
pour certains $z_i$ dans $A$. De plus, pour $0 < i < n$, on a
$$
\text{d}(\gamma_i(y) z_i) =
\text{d}(\gamma_i(y))z_i + \gamma_i(y)\text{d}(z_i) =
\text{d}(y) \gamma_{i - 1}(y) z_i + \gamma_i(y)\text{d}(z_i)
$$
et le premier terme a une image nulle dans $B$, car $\text{d}(y)$ a une
image nulle dans $B$. Ainsi, $\gamma_n(x')$ et $\gamma_n(x)$ ont la même
image dans $H(B)$. On obtient donc une application bien définie
$\gamma_n : H_{2d}(B) \to H_{2nd}(B)$ pour tous $d > 0$ et $n > 0$.
Nous omettons la vérification du fait qu'elle définit une structure de
puissances divisées sur $H(B)$.
\end{proof}

\begin{lemma}
\label{lemma-base-change-div}
Soit $(A, \text{d}, \gamma)$ comme dans la définition
\ref{definition-divided-powers-dga}. Soit $R \to R'$ un homomorphisme
d'anneaux. Alors $\text{d}$ et $\gamma$ induisent des structures analogues
sur $A' = A \otimes_R R'$ telles que $(A', \text{d}, \gamma)$ soit comme
dans la définition \ref{definition-divided-powers-dga}.
\end{lemma}

\begin{proof}
Remarquons que $A'_{even} = A_{even} \otimes_R R'$ et
$A'_{even, +} = A_{even, +} \otimes_R R'$. Il s'agit donc de montrer que
les puissances divisées $\gamma$ se prolongent à $A'_{even}$ (au sens de la
définition \ref{definition-extends}). Une fois établi que $\gamma$ se
prolonge, il est facile de voir que ce prolongement possède toutes les
propriétés voulues.

\medskip\noindent
Choisissons une $R$-algèbre polynomiale $P$ sur un ensemble quelconque de
générateurs, ainsi qu'une surjection de $R$-algèbres
$P \to R'$. L'homomorphisme d'anneaux
$A_{even} \to A_{even} \otimes_R P$ est plat ; les puissances divisées
$\gamma$ se prolongent donc de manière unique à
$A_{even} \otimes_R P$ par le lemme \ref{lemma-gamma-extends}.
Soit $J = \Ker(P \to R')$. Pour montrer que $\gamma$ se prolonge à
$A \otimes_R R'$, il suffit de montrer que
$I' = \Ker(A_{even, +} \otimes_R P \to A_{even, +} \otimes_R R')$
est engendré par des éléments $z$ tels que $\gamma_n(z) \in I'$ pour tout
$n > 0$. C'est clair, car $I'$ est engendré par les éléments de la forme
$x \otimes f$, avec $x \in A_{even, +}$ et $f \in \Ker(P \to R')$.
\end{proof}

\begin{lemma}
\label{lemma-extend-differential}
Soit $(A, \text{d}, \gamma)$ comme dans la définition
\ref{definition-divided-powers-dga}. Soit $d \geq 1$ un entier. Soit
$A\langle T \rangle$ l'algèbre polynomiale graduée à puissances divisées en
$T$, avec $\deg(T) = d$, construite dans l'exemple
\ref{example-adjoining-odd} ou \ref{example-adjoining-even}.
Soit $f \in A_{d - 1}$ un élément tel que $\text{d}(f) = 0$. Il existe une
unique différentielle $\text{d}$ sur $A\langle T\rangle$ telle que
$\text{d}(T) = f$ et telle que $\text{d}$ soit compatible avec la structure
de puissances divisées sur $A\langle T \rangle$.
\end{lemma}

\begin{proof}
Cela résulte d'un calcul direct que nous omettons.
\end{proof}

\noindent
Dans le lemme \ref{lemma-compute-cohomology-adjoin-variable}, nous calculerons
l'homologie de $A\langle T \rangle$ dans certains cas particuliers. Voici la
construction de Tate, telle qu'elle a été étendue par Avramov et Halperin.

\begin{lemma}
\label{lemma-tate-resolution}
Soit $R \to S$ un homomorphisme d'anneaux commutatifs. Il existe une
factorisation
$$
R \to A \to S
$$
ayant les propriétés suivantes :
\begin{enumerate}
\item $(A, \text{d}, \gamma)$ est comme dans la définition
\ref{definition-divided-powers-dga},
\item $A \to S$ est un quasi-isomorphisme (si l'on munit $S$ de la
différentielle nulle),
\item $A_0 = R[x_j: j\in J] \to S$ est une surjection quelconque d'un anneau
de polynômes sur $S$, et
\item $A$ est une algèbre polynomiale graduée à puissances divisées sur $R$.
\end{enumerate}
La dernière condition signifie que $A$ est construite à partir de $A_0$ en
adjoignant successivement un ensemble de variables $T$ dans chaque degré
$> 0$, comme dans l'exemple \ref{example-adjoining-odd} ou
\ref{example-adjoining-even}. De plus, si $R$ est noethérien et si
$R\to S$ est de type fini, on peut choisir $A$ avec seulement un nombre fini
de générateurs dans chaque degré.
\end{lemma}

\begin{proof}
Détaillons la construction lorsque $R$ est noethérien et $R\to S$ est de
type fini. Sans ces hypothèses, la démonstration est la même, sauf qu'il faut
utiliser dans chaque degré un ensemble de générateurs éventuellement infini.

\medskip\noindent
Début de la construction : soit $A(0) = R[x_1, \ldots, x_n]$ un anneau de
polynômes usuel, et soit $A(0) \to S$ une surjection. Pour la graduation, on
pose $A(0)_0 = A(0)$ et $A(0)_d = 0$ pour $d \not = 0$. Ainsi,
$\text{d} = 0$ et $\gamma_n$ est également nulle pour $n > 0$.

\medskip\noindent
Choisissons des générateurs $f_1, \ldots, f_m \in R[x_1, \ldots, x_n]$ du
noyau de l'homomorphisme donné
$A(0) = R[x_1, \ldots, x_n] \to S$. En appliquant l'exemple
\ref{example-adjoining-odd} $m$ fois, on obtient
$$
A(1) = A(0)\langle T_1, \ldots, T_m\rangle
$$
avec $\deg(T_i) = 1$, en tant qu'algèbre polynomiale graduée à puissances
divisées. Posons $\text{d}(T_i) = f_i$. Comme $A(1)$ est une algèbre
polynomiale à puissances divisées sur $A(0)$ et que
$\text{d}(f_i) = 0$, ceci se prolonge de manière unique en une différentielle
sur $A(1)$, d'après le lemme \ref{lemma-extend-differential}.

\medskip\noindent
Hypothèse de récurrence : supposons données des factorisations
$$
R \to A(0) \to A(1) \to \ldots \to A(m) \to S
$$
où $A(0)$ et $A(1)$ sont comme ci-dessus, et telles que chaque
$R \to A(m') \to S$, pour $2 \leq m' \leq m$, vérifie les propriétés (1)
et (4) de l'énoncé du lemme, la propriété (2) étant remplacée par la condition
que $H_i(A(m')) \to H_i(S)$ est un isomorphisme pour
$m' > i \geq 0$. Le cas initial est $m = 1$.

\medskip\noindent
Étape de récurrence : supposons donné $R \to A(m) \to S$ comme dans
l'hypothèse de récurrence. Considérons le groupe $H_m(A(m))$. C'est un module
sur $H_0(A(m)) = S$. En fait, c'est un sous-quotient de $A(m)_m$, qui est un
module de type fini sur $A(m)_0 = R[x_1, \ldots, x_n]$. On peut donc choisir
un nombre fini d'éléments
$$
e_1, \ldots, e_t \in \Ker(\text{d} : A(m)_m \to A(m)_{m - 1})
$$
dont les images engendrent ce module. En appliquant l'exemple
\ref{example-adjoining-odd} ou \ref{example-adjoining-even} $t$ fois, on
obtient
$$
A(m + 1) = A(m)\langle T_1, \ldots, T_t\rangle
$$
avec $\deg(T_i) = m + 1$, en tant qu'algèbre graduée à puissances divisées.
Posons $\text{d}(T_i) = e_i$. Comme $A(m+1)$ est une algèbre polynomiale à
puissances divisées sur $A(m)$ et que $\text{d}(e_i) = 0$, ceci se prolonge
de manière unique en une différentielle sur $A(m + 1)$ compatible avec la
structure de puissances divisées. Puisque nous n'avons ajouté des éléments
qu'en degré $m + 1$ et au-dessus, on a
$H_i(A(m + 1)) = H_i(A(m))$ pour $i < m$. De plus,
$H_m(A(m + 1)) = 0$ par construction.

\medskip\noindent
Pour achever la démonstration, remarquons que nous avons construit une suite
de morphismes
$$
R \to A(0) \to A(1) \to \ldots \to A(m) \to A(m + 1) \to \ldots \to S
$$
et posons $A = \colim A(m)$.
\end{proof}

\begin{lemma}
\label{lemma-tate-resoluton-pseudo-coherent-ring-map}
Soit $R \to S$ un homomorphisme pseudo-cohérent d'anneaux (Compléments sur
l'algèbre, définition
\ref{more-algebra-definition-pseudo-coherent-perfect}). Alors le
lemme \ref{lemma-tate-resolution} vaut avec une résolution $A$ de $S$
possédant un nombre fini de générateurs dans chaque degré.
\end{lemma}

\begin{proof}
La démonstration est exactement la même que celle du
lemme \ref{lemma-tate-resolution}. Le seul point supplémentaire est que,
étant donné $A(m) \to S$, il faut montrer que
$H_m = H_m(A(m))$ est un $R[x_1, \ldots, x_m]$-module de type fini (de sorte
qu'à l'étape suivante il ne soit nécessaire d'ajouter qu'un nombre fini de
variables). Considérons le complexe
$$
\ldots \to A(m)_{m - 1} \to A(m)_m \to A(m)_{m - 1} \to
\ldots \to A(m)_0 \to S \to 0
$$
Comme $S$ est un $R[x_1, \ldots, x_n]$-module pseudo-cohérent et que
$A(m)_i$ est un $R[x_1, \ldots, x_n]$-module libre de type fini, ce complexe
est pseudo-cohérent ; voir Compléments sur l'algèbre, lemme
\ref{more-algebra-lemma-complex-pseudo-coherent-modules}. Comme le complexe
est exact en degrés homologiques $> m$, on conclut que $H_m$ est un
$R$-module de type fini, d'après Compléments sur l'algèbre, lemme
\ref{more-algebra-lemma-finite-cohomology}.
\end{proof}

\begin{lemma}
\label{lemma-uniqueness-tate-resolution}
Soit $R$ un anneau commutatif. Supposons que $(A, \text{d}, \gamma)$ et
$(B, \text{d}, \gamma)$ soient comme dans la définition
\ref{definition-divided-powers-dga}. Soit
$\overline{\varphi} : H_0(A) \to H_0(B)$ un homomorphisme de $R$-algèbres.
Supposons que
\begin{enumerate}
\item $A$ soit une algèbre polynomiale graduée à puissances divisées sur $R$.
\item $H_k(B) = 0$ pour $k > 0$.
\end{enumerate}
Il existe alors un morphisme $\varphi : A \to B$ de $R$-algèbres
différentielles graduées, compatible avec les puissances divisées, qui relève
$\overline{\varphi}$.
\end{lemma}

\begin{proof}
L'hypothèse signifie que $A$ s'obtient à partir de $R$ en adjoignant
successivement un ensemble de générateurs polynomiaux en degré zéro, des
générateurs extérieurs en degrés impairs positifs et des générateurs à
puissances divisées en degrés pairs positifs. On dispose donc d'une filtration
$R \subset A(0) \subset A(1) \subset \ldots$
de $A$ telle que $A(m + 1)$ s'obtienne à partir de $A(m)$ en adjoignant des
générateurs du type approprié (que nous appellerons simplement
« générateurs à puissances divisées ») en degré $m + 1$. En particulier,
$A(0) \to H_0(A)$ est un morphisme surjectif d'une algèbre polynomiale
usuelle sur $R$ vers $H_0(A)$. On peut donc relever $\overline{\varphi}$ en un
homomorphisme de $R$-algèbres $\varphi(0) : A(0) \to B_0$.

\medskip\noindent
Écrivons $A(1) = A(0)\langle T_j:j\in J\rangle$ pour un ensemble $J$ de
variables à puissances divisées $T_j$ de degré $1$. Soit $f_j \in B_0$
défini par $f_j = \varphi(0)(\text{d}(T_j))$. Remarquons que l'image de $f_j$
dans $H_0(B)$ est nulle, puisque l'image de $\text{d}T_j$ dans $H_0(A)$ est
nulle. Il existe donc $b_j \in B_1$ tel que $\text{d}(b_j) = f_j$.
D'après la propriété universelle des algèbres polynomiales à puissances
divisées du lemme \ref{lemma-divided-power-polynomial-algebra}, il existe un
relèvement $\varphi(1) : A(1) \to B$ de $\varphi(0)$ qui envoie $T_j$ sur
$f_j$.

\medskip\noindent
Après avoir construit $\varphi(m)$ pour un certain $m \geq 1$, on construit
$\varphi(m + 1) : A(m + 1) \to B$ exactement de la même manière. Nous
omettons les détails.
\end{proof}

\begin{lemma}
\label{lemma-divided-powers-on-tor}
Soit $R$ un anneau commutatif. Soient $S$ et $T$ des $R$-algèbres
commutatives. Il existe alors une structure canonique de $R$-algèbre
strictement commutative au sens gradué et munie de puissances divisées sur
$$
\operatorname{Tor}_*^R(S, T).
$$
\end{lemma}

\begin{proof}
Choisissons une factorisation $R \to A \to S$ comme ci-dessus. Comme
$A \to S$ est un quasi-isomorphisme et que $A_d$ est un $R$-module libre,
l'algèbre différentielle graduée $B = A \otimes_R T$ calcule les groupes de
Tor affichés dans le lemme. Choisissons une surjection
$R[y_j:j\in J] \to T$. Alors $B$ est un quotient de l'algèbre différentielle
graduée $A[y_j:j\in J]$, dont l'homologie est concentrée en degré $0$ (elle
est égale à $S[y_j:j\in J]$). D'après le
lemme \ref{lemma-base-change-div}, les algèbres différentielles graduées $B$
et $A[y_j:j\in J]$ possèdent des structures de puissances divisées
compatibles avec les différentielles. On obtient donc la structure de
puissances divisées cherchée sur $H(B)$ par le
lemme \ref{lemma-dpdga-good}.

\medskip\noindent
La structure d'algèbre à puissances divisées ainsi construite ne dépend pas
du choix de $A$. En effet, si $A'$ est un second choix, le
lemme \ref{lemma-uniqueness-tate-resolution} fournit un morphisme
$A \to A'$ qui préserve toutes les structures ainsi que les augmentations
vers $S$. Le morphisme induit
$B = A \otimes_R T \to A' \otimes_R T' = B'$ préserve alors lui aussi toutes
les structures et est un quasi-isomorphisme. L'isomorphisme induit entre les
algèbres de Tor est donc compatible avec les produits et les puissances
divisées.
\end{proof}
\section{Application aux intersections complètes}
\label{section-application-ci}

\noindent
Soit $R$ un anneau. Soit $(A, \text{d}, \gamma)$ comme dans la
définition \ref{definition-divided-powers-dga}.
Une {\it dérivation} de degré $2$ est une application $R$-linéaire
$\theta : A \to A$ possédant les propriétés suivantes :
\begin{enumerate}
\item $\theta(A_d) \subset A_{d - 2}$,
\item $\theta(xy) = \theta(x)y + x\theta(y)$,
\item $\theta$ commute avec $\text{d}$,
\item $\theta(\gamma_n(x)) = \theta(x) \gamma_{n - 1}(x)$
pour tout $x \in A_{2d}$ et tout $d$.
\end{enumerate}
Dans le lemme suivant, nous construisons une dérivation.

\begin{lemma}
\label{lemma-get-derivation}
Soit $R$ un anneau. Soit $(A, \text{d}, \gamma)$ comme dans la
définition \ref{definition-divided-powers-dga}.
Soit $R' \to R$ une surjection d'anneaux dont le noyau
est de carré nul et engendré par un élément $f$.
Si $A$ est une algèbre polynomiale graduée à puissances divisées sur $R$,
ayant un nombre fini de variables en chaque degré,
alors on obtient une dérivation
$\theta : A/IA \to A/IA$, où $I$ est l'annulateur
de $f$ dans $R$.
\end{lemma}

\begin{proof}
Comme $A$ est une algèbre polynomiale à puissances divisées, il existe une
algèbre polynomiale à puissances divisées $A'$ sur $R'$ telle que
$A = A' \otimes_{R'} R$.
De plus, on peut relever $\text{d}$ en un opérateur
$R$-linéaire $\text{d}$ sur $A'$ tel que
\begin{enumerate}
\item  $\text{d}(xy) = \text{d}(x)y + (-1)^{\deg(x)}x \text{d}(y)$
pour $x, y \in A'$ homogènes, et
\item  $\text{d}(\gamma_n(x)) = \text{d}(x) \gamma_{n - 1}(x)$ pour
$x \in A'_{even, +}$.
\end{enumerate}
Nous omettons les détails (indication : procéder une variable à la fois).
Il se peut toutefois que $\text{d}^2$
ne soit pas nul sur $A'$. Il est clair que $\text{d}^2$ envoie $A'$ dans
$fA' \cong A/IA$. Par conséquent, $\text{d}^2$ annule $fA'$ et se factorise
en une application $A \to A/IA$. Comme $\text{d}^2$ est $R$-linéaire, on
obtient l'application voulue $\theta : A/IA \to A/IA$. La vérification des
propriétés d'une dérivation est immédiate.
\end{proof}

\begin{lemma}
\label{lemma-compute-theta}
On conserve les hypothèses et les notations du lemme
\ref{lemma-get-derivation}.
Supposons que $S = H_0(A)$ soit isomorphe à
$R[x_1, \ldots, x_n]/(f_1, \ldots, f_m)$
pour certains $n$, $m$ et $f_j \in R[x_1, \ldots, x_n]$.
Supposons en outre donnée une relation
$$
\sum r_j f_j = 0
$$
avec $r_j \in R[x_1, \ldots, x_n]$.
Choisissons des relèvements $r'_j, f'_j \in R'[x_1, \ldots, x_n]$ de
$r_j, f_j$. Écrivons $\sum r'_j f'_j = gf$ pour un certain
$g \in R/I[x_1, \ldots, x_n]$.
Si $H_1(A) = 0$ et si tous les coefficients de chaque $r_j$ appartiennent
à $I$, alors il existe un élément $\xi \in H_2(A/IA)$ tel que
$\theta(\xi) = g$ dans $S/IS$.
\end{lemma}

\begin{proof}
Soit $A(0) \subset A(1) \subset A(2) \subset \ldots$ la filtration
de $A$ telle que $A(m)$ s'obtienne à partir de $A(m - 1)$ en adjoignant
des variables à puissances divisées de degré $m$. Alors $A(0)$ est une
algèbre polynomiale sur $R$, munie d'une surjection de $R$-algèbres
$A(0) \to S$.
On peut donc choisir une application
$$
\varphi : R[x_1, \ldots, x_n] \to A(0)
$$
qui relève les applications d'augmentation vers $S$.
Ensuite, $A(1) = A(0)\langle T_1, \ldots, T_t \rangle$, où les $T_i$
sont des variables à puissances divisées de degré $1$. Comme
$H_0(A) = S$, on peut choisir des éléments
$\xi_j \in \sum A(0)T_i$ tels que $\text{d}(\xi_j) = \varphi(f_j)$.
Alors
$$
\text{d}\left(\sum \varphi(r_j) \xi_j\right) =
\sum  \varphi(r_j) \varphi(f_j) =  \sum \varphi(r_jf_j) = 0
$$
Comme $H_1(A) = 0$, on peut choisir $\xi \in A_2$ tel que
$\text{d}(\xi) = \sum \varphi(r_j) \xi_j$.
Si les coefficients de $r_j$ appartiennent à $I$, alors
$\varphi(r_j)$ possède aussi cette propriété. Dans ce cas,
$\text{d}(\xi)$ s'annule dans $A_1/IA_1$ et
$\xi$ définit donc une classe dans $H_2(A/IA)$.

\medskip\noindent
La construction de $\theta$ dans la démonstration du lemme
\ref{lemma-get-derivation}
consiste à relever successivement $A(i)$ en $A'(i)$ ainsi qu'à relever
la différentielle $\text{d}$. On relève $\varphi$
en $\varphi' : R'[x_1, \ldots, x_n] \to A'(0)$.
Ensuite, on a $A'(1) = A'(0)\langle T_1, \ldots, T_t\rangle$.
De plus, on peut relever $\xi_j$ en
$\xi'_j \in \sum A'(0)T_i$.
Alors $\text{d}(\xi'_j) = \varphi'(f'_j) + f a_j$ pour un certain
$a_j \in A'(0)$.
Considérons un relèvement $\xi' \in A'_2$ de $\xi$.
On sait alors que
$$
\text{d}(\xi') = \sum \varphi'(r'_j)\xi'_j + \sum fb_iT_i
$$
pour certains $b_i \in A(0)$. En appliquant de nouveau $\text{d}$,
on obtient
$$
\theta(\xi) = \sum \varphi'(r'_j)\varphi'(f'_j) +
\sum f \varphi'(r'_j) a_j + \sum fb_i \text{d}(T_i)
$$
Le premier terme donne le résultat voulu. Le deuxième est nul, car les
coefficients de $r_j$ appartiennent à $I$ et sont donc annulés par $f$.
Le troisième a une image nulle dans $H_0$, car l'image de
$\text{d}(T_i)$ y est nulle.
\end{proof}

\noindent
La méthode de démonstration du lemme suivant semble due à Gulliksen.

\begin{lemma}
\label{lemma-not-finite-pd}
Soit $R' \to R$ une surjection d'anneaux noethériens dont le noyau est de
carré nul et engendré par un élément $f$. Soit
$S = R[x_1, \ldots, x_n]/(f_1, \ldots, f_m)$.
Soit $\sum r_j f_j = 0$ une relation dans $R[x_1, \ldots, x_n]$.
Supposons que
\begin{enumerate}
\item chaque $r_j$ ait ses coefficients dans l'annulateur $I$ de $f$ dans $R$,
\item pour certains relèvements $r'_j, f'_j \in R'[x_1, \ldots, x_n]$, on ait
$\sum r'_j f'_j = gf$, où $g$ n'est pas nilpotent dans $S/IS$.
\end{enumerate}
Alors $S$ n'est pas de dimension de Tor finie sur $R$ (autrement dit,
$S$ n'est pas une $R$-algèbre parfaite).
\end{lemma}

\begin{proof}
Choisissons une résolution de Tate $R \to A \to S$ comme dans le
lemme \ref{lemma-tate-resolution}.
Soient $\xi \in H_2(A/IA)$ et $\theta : A/IA \to A/IA$ l'élément
et la dérivation obtenus dans les lemmes \ref{lemma-get-derivation} et
\ref{lemma-compute-theta}.
On observe que
$$
\theta^n(\gamma_n(\xi)) = g^n
$$
dans $H_0(A/IA) = S/IS$.
Ainsi, si $g$ n'est pas nilpotent dans $S/IS$, alors $\xi^n$ est non nul dans
$H_{2n}(A/IA)$ pour tout $n > 0$. Comme
$H_{2n}(A/IA) = \text{Tor}^R_{2n}(S, R/I)$, on conclut.
\end{proof}

\noindent
Le résultat suivant se trouve dans \cite{Rodicio}.

\begin{lemma}
\label{lemma-injective}
Soit $(A, \mathfrak m)$ un anneau local noethérien. Soient
$I \subset J \subset A$ des idéaux propres. Si $A/J$ est de dimension
de Tor finie sur $A/I$, alors l'application
$I/\mathfrak m I \to J/\mathfrak m J$ est injective.
\end{lemma}

\begin{proof}
Soit $f \in I$ un élément dont l'image dans $I/\mathfrak m I$ est non nulle
et dont l'image dans $J/\mathfrak mJ$ est nulle. On peut choisir un idéal $I'$
tel que $\mathfrak mI \subset I' \subset I$ et que $I/I'$ soit engendré par
l'image de $f$. Posons $R = A/I$ et $R' = A/I'$. Écrivons
$J = (a_1, \ldots, a_m)$
avec $a_j \in A$. Alors $f = \sum b_j a_j$ pour certains
$b_j \in \mathfrak m$.
Soient $r_j, f_j \in R$ resp.\ $r'_j, f'_j \in R'$ les images de
$b_j, a_j$. On se trouve alors
dans la situation du lemme \ref{lemma-not-finite-pd}
(l'idéal $I$ de ce lemme étant égal à $\mathfrak m_R$),
ce qui démontre le lemme.
\end{proof}

\begin{lemma}
\label{lemma-regular-sequence}
Soit $(A, \mathfrak m)$ un anneau local noethérien. Soient
$I \subset J \subset A$ des idéaux propres. Supposons que
\begin{enumerate}
\item $A/J$ soit de dimension de Tor finie sur $A/I$, et
\item $J$ soit engendré par une suite régulière.
\end{enumerate}
Alors $I$ est engendré par une suite régulière et $J/I$
est engendré par une suite régulière.
\end{lemma}

\begin{proof}
Le lemme \ref{lemma-injective} montre que l'application
$I/\mathfrak m I \to J/\mathfrak m J$
est injective. Il existe donc $s \leq r$ et un système minimal
de générateurs $f_1, \ldots, f_r$ de $J$ tels que
$f_1, \ldots, f_s$ appartiennent à $I$
et forment un système minimal de générateurs de $I$.
Le lemme en résulte, car tout système minimal de générateurs de $J$
est une suite régulière d'après
Compléments sur l'algèbre, lemmes
\ref{more-algebra-lemma-independence-of-generators} et
\ref{more-algebra-lemma-noetherian-finite-all-equivalent}.
\end{proof}

\begin{lemma}
\label{lemma-perfect-map-ci}
Soit $R \to S$ un homomorphisme local d'anneaux locaux noethériens.
Soient $I \subset R$ et $J \subset S$ des idéaux tels que
$IS \subset J$. Si $R \to S$ est plat et si $S/\mathfrak m_RS$ est
régulier, alors les propriétés suivantes sont équivalentes :
\begin{enumerate}
\item $J$ est engendré par une suite régulière et
$S/J$ est de dimension de Tor finie comme module sur $R/I$,
\item $J$ est engendré par une suite régulière et
$\text{Tor}^{R/I}_p(S/J, R/\mathfrak m_R)$ n'est non nul
que pour un nombre fini d'entiers $p$,
\item $I$ est engendré par une suite régulière
et $J/IS$ est engendré par une suite régulière dans $S/IS$.
\end{enumerate}
\end{lemma}

\begin{proof}
Si (3) est vérifiée, alors $J$ est engendré par une suite régulière ; voir,
par exemple, Compléments sur l'algèbre, lemmes
\ref{more-algebra-lemma-join-koszul-regular-sequences} et
\ref{more-algebra-lemma-noetherian-finite-all-equivalent}.
De plus, sous cette même hypothèse, $S/J = (S/I)/(J/I)$
est de dimension projective finie sur $S/IS$, car le complexe de Koszul
fournit une résolution libre finie de $S/J$ sur $S/IS$.
Comme $R/I \to S/IS$ est plat, il s'ensuit que $S/J$ est de dimension
de Tor finie sur $R/I$, d'après
Compléments sur l'algèbre, lemme
\ref{more-algebra-lemma-flat-push-tor-amplitude}.
Ainsi, (3) implique (1).

\medskip\noindent
L'implication (1) $\Rightarrow$ (2) est immédiate.
Supposons (2). D'après
Compléments sur l'algèbre, lemme
\ref{more-algebra-lemma-perfect-over-regular-local-ring},
$S/J$ est de dimension de Tor finie sur $S/IS$.
On peut donc appliquer le lemme \ref{lemma-regular-sequence}
pour conclure que $IS$ et $J/IS$ sont engendrés par des suites régulières.
Soit $f_1, \ldots, f_r \in I$ un système minimal de générateurs de $I$.
Comme $R \to S$ est plat, les éléments $f_1, \ldots, f_r$ forment un système
minimal de générateurs de $IS$ dans $S$. Ainsi, $f_1, \ldots, f_r \in R$
est une suite d'éléments dont les images dans $S$ forment une suite régulière,
d'après Compléments sur l'algèbre, lemmes
\ref{more-algebra-lemma-independence-of-generators} et
\ref{more-algebra-lemma-noetherian-finite-all-equivalent}.
Par conséquent, $f_1, \ldots, f_r$ est une suite régulière dans $R$,
d'après Algèbre, lemme \ref{algebra-lemma-flat-increases-depth}.
\end{proof}





\section{Anneaux d'intersection complète locale}
\label{section-lci}

\noindent
Soit $(A, \mathfrak m)$ un anneau local noethérien complet.
D'après le théorème de structure de Cohen (voir
Algèbre, théorème \ref{algebra-theorem-cohen-structure-theorem}),
on peut écrire $A$ comme quotient d'un anneau local noethérien
régulier et complet $R$. On dit que $A$ est une
{\it intersection complète}
s'il existe une surjection $R \to A$,
où $R$ est un anneau local régulier, dont le noyau
est engendré par une suite régulière.
Le lemme suivant montre que cette notion ne dépend pas
du choix de la surjection.

\begin{lemma}
\label{lemma-ci-well-defined}
Soit $(A, \mathfrak m)$ un anneau local noethérien complet.
Les propriétés suivantes sont équivalentes :
\begin{enumerate}
\item pour toute surjection d'anneaux locaux $R \to A$, où $R$
est un anneau local régulier, le noyau de $R \to A$ est engendré
par une suite régulière, et
\item pour une certaine surjection d'anneaux locaux $R \to A$, où $R$
est un anneau local régulier, le noyau de $R \to A$ est engendré
par une suite régulière.
\end{enumerate}
\end{lemma}

\begin{proof}
Soit $k$ le corps résiduel de $A$. Si la caractéristique de
$k$ est $p > 0$, on désigne par $\Lambda$ un anneau de Cohen
(Algèbre, définition \ref{algebra-definition-cohen-ring})
de corps résiduel $k$ (Algèbre, lemme
\ref{algebra-lemma-cohen-rings-exist}).
Si la caractéristique de $k$ est $0$, on pose $\Lambda = k$.
Rappelons que $\Lambda[[x_1, \ldots, x_n]]$, pour tout $n$,
est formellement lisse sur $\mathbf{Z}$, resp.\ $\mathbf{Q}$,
pour la topologie $\mathfrak m$-adique ; voir
Compléments sur l'algèbre, lemme
\ref{more-algebra-lemma-power-series-ring-over-Cohen-fs}.
Fixons une surjection $\Lambda[[x_1, \ldots, x_n]] \to A$ donnée par
le théorème de structure de Cohen
(Algèbre, théorème \ref{algebra-theorem-cohen-structure-theorem}).

\medskip\noindent
Soit $R \to A$ une surjection, où $R$ est un anneau local régulier.
Soit $f_1, \ldots, f_r$ un système minimal de générateurs
de $\Ker(R \to A)$. Nous utiliserons sans autre mention
le fait qu'un idéal d'un anneau local noethérien est engendré par une suite
régulière si et seulement si tout système minimal de générateurs est une
suite régulière. On observe que $f_1, \ldots, f_r$
est une suite régulière dans $R$ si et seulement si $f_1, \ldots, f_r$
est une suite régulière dans le complété $R^\wedge$, d'après
Algèbre, lemmes \ref{algebra-lemma-flat-increases-depth} et
\ref{algebra-lemma-completion-flat}.
De plus,
$$
R^\wedge/(f_1, \ldots, f_r)R^\wedge =
(R/(f_1, \ldots, f_n))^\wedge = A^\wedge = A
$$
car $A$ est complet pour la topologie $\mathfrak m_A$-adique
(la première égalité résulte d'Algèbre, lemme
\ref{algebra-lemma-completion-tensor}). Enfin,
l'anneau $R^\wedge$ est régulier puisque $R$ est régulier
(Compléments sur l'algèbre, lemme
\ref{more-algebra-lemma-completion-regular}).
On peut donc supposer que $R$ est complet.

\medskip\noindent
Si $R$ est complet, on peut choisir une application
$\Lambda[[x_1, \ldots, x_n]] \to R$ relevant l'application donnée
$\Lambda[[x_1, \ldots, x_n]] \to A$ ; voir
Compléments sur l'algèbre, lemme \ref{more-algebra-lemma-lift-continuous}.
En ajoutant des variables $y_1, \ldots, y_m$ envoyées
sur des générateurs du noyau de $R \to A$, on peut supposer que
$\Lambda[[x_1, \ldots, x_n, y_1, \ldots, y_m]] \to R$ est surjective
(nous omettons quelques détails). On peut alors considérer le diagramme
commutatif
$$
\xymatrix{
\Lambda[[x_1, \ldots, x_n, y_1, \ldots, y_m]] \ar[r] \ar[d] & R \ar[d] \\
\Lambda[[x_1, \ldots, x_n]] \ar[r] & A
}
$$
Le lemme \ref{algebra-lemma-ci-well-defined} d'Algèbre montre que
la condition relative à $R \to A$ équivaut à celle relative à
l'application fixée
$\Lambda[[x_1, \ldots, x_n]] \to A$. Cela achève la démonstration du lemme.
\end{proof}

\noindent
Les deux lemmes suivants constituent des vérifications de cohérence
de la définition donnée ci-dessus.

\begin{lemma}
\label{lemma-quotient-regular-ring-by-regular-sequence}
Soit $R$ un anneau régulier. Soit $\mathfrak p \subset R$ un idéal premier.
Soit $f_1, \ldots, f_r \in \mathfrak p$ une suite régulière.
Alors le complété de
$$
A = (R/(f_1, \ldots, f_r))_\mathfrak p =
R_\mathfrak p/(f_1, \ldots, f_r)R_\mathfrak p
$$
est une intersection complète au sens défini ci-dessus.
\end{lemma}

\begin{proof}
Le complété de $A$ est égal à
$A^\wedge = R_\mathfrak p^\wedge/(f_1, \ldots, f_r)R_\mathfrak p^\wedge$,
car le foncteur de complétion est exact sur les modules finis sur l'anneau
noethérien $R_\mathfrak p$
(Algèbre, lemme \ref{algebra-lemma-completion-tensor}).
L'image de la suite $f_1, \ldots, f_r$ dans $R_\mathfrak p$
est une suite régulière d'après
Algèbre, lemmes \ref{algebra-lemma-completion-flat} et
\ref{algebra-lemma-flat-increases-depth}.
De plus, $R_\mathfrak p^\wedge$ est un anneau local régulier d'après
Compléments sur l'algèbre, lemme
\ref{more-algebra-lemma-completion-regular}.
Le résultat découle donc de notre définition d'une intersection
complète pour les anneaux locaux complets.
\end{proof}

\noindent
Le lemme suivant est l'analogue d'Algèbre, lemme
\ref{algebra-lemma-lci}.

\begin{lemma}
\label{lemma-quotient-regular-ring}
Soit $R$ un anneau régulier. Soit $\mathfrak p \subset R$ un idéal premier.
Soit $I \subset \mathfrak p$ un idéal.
Posons $A = (R/I)_\mathfrak p = R_\mathfrak p/I_\mathfrak p$.
Les propriétés suivantes sont équivalentes :
\begin{enumerate}
\item le complété de $A$
est une intersection complète au sens ci-dessus,
\item $I_\mathfrak p \subset R_\mathfrak p$ est engendré
par une suite régulière,
\item le module $(I/I^2)_\mathfrak p$ peut être engendré par
$\dim(R_\mathfrak p) - \dim(A)$ éléments,
\item ajouter quelque chose ici.
\end{enumerate}
\end{lemma}

\begin{proof}
On peut remplacer, et l'on remplace, $R$ par son localisé en
$\mathfrak p$.
Alors $\mathfrak p = \mathfrak m$ est l'idéal maximal de $R$
et $A = R/I$. Soit $f_1, \ldots, f_r \in I$ un système minimal
de générateurs. Le complété de $A$ est égal à
$A^\wedge = R^\wedge/(f_1, \ldots, f_r)R^\wedge$,
car le foncteur de complétion est exact sur les modules finis sur l'anneau
noethérien $R_\mathfrak p$
(Algèbre, lemme \ref{algebra-lemma-completion-tensor}).

\medskip\noindent
Si (1) est vérifiée, alors l'image de la suite $f_1, \ldots, f_r$ dans
$R^\wedge$ est une suite régulière par hypothèse. C'est donc une suite
régulière dans $R$ d'après Algèbre, lemmes
\ref{algebra-lemma-completion-flat} et
\ref{algebra-lemma-flat-increases-depth}. Ainsi, (1) implique (2).

\medskip\noindent
Supposons (3). Posons $c = \dim(R) - \dim(A)$ et choisissons
$f_1, \ldots, f_c \in I$ dont les images engendrent $I/I^2$.
Le lemme de Nakayama
(Algèbre, lemme \ref{algebra-lemma-NAK})
montre que $I = (f_1, \ldots, f_c)$. Comme $R$ est régulier, donc
de Cohen--Macaulay (Algèbre, proposition
\ref{algebra-proposition-CM-module}),
$f_1, \ldots, f_c$ est une suite régulière d'après
Algèbre, proposition \ref{algebra-proposition-CM-module}.
Ainsi, (3) implique (2).
Enfin, (2) implique (1) d'après le
lemme \ref{lemma-quotient-regular-ring-by-regular-sequence}.
\end{proof}

\noindent
Le résultat suivant est dû à Avramov ; voir \cite{Avramov}.

\begin{proposition}
\label{proposition-avramov}
Soit $A \to B$ un homomorphisme local plat d'anneaux locaux noethériens.
Les propriétés suivantes sont équivalentes :
\begin{enumerate}
\item $B^\wedge$ est une intersection complète,
\item $A^\wedge$ et $(B/\mathfrak m_A B)^\wedge$ sont des intersections complètes.
\end{enumerate}
\end{proposition}

\begin{proof}
Considérons le diagramme
$$
\xymatrix{
B \ar[r] & B^\wedge \\
A \ar[u] \ar[r] & A^\wedge \ar[u]
}
$$
Comme les applications horizontales sont fidèlement plates
(Algèbre, lemme \ref{algebra-lemma-completion-faithfully-flat}),
la flèche verticale de droite est plate
(par exemple, d'après Algèbre, lemme
\ref{algebra-lemma-criterion-flatness-fibre-Noetherian}).
De plus,
$(B/\mathfrak m_A B)^\wedge = B^\wedge/\mathfrak m_{A^\wedge} B^\wedge$
d'après Algèbre, lemme \ref{algebra-lemma-completion-tensor}.
On peut donc supposer que $A$ et $B$ sont des anneaux locaux noethériens
complets.

\medskip\noindent
Supposons que $A$ et $B$ soient des anneaux locaux noethériens complets.
Choisissons un diagramme
$$
\xymatrix{
S \ar[r] & B \\
R \ar[u] \ar[r] & A \ar[u]
}
$$
comme dans Compléments sur l'algèbre, lemme
\ref{more-algebra-lemma-embed-map-Noetherian-complete-local-rings}.
Posons $I = \Ker(R \to A)$ et $J = \Ker(S \to B)$.
Comme $R/I = A \to B = S/J$ est plat, l'application
$J/IS \otimes_R R/\mathfrak m_R \to J/J \cap \mathfrak m_R S$
est un isomorphisme. Par conséquent, un système minimal de générateurs
de $J/IS$ s'envoie sur un système minimal de générateurs de
$\Ker(S/\mathfrak m_R S \to B/\mathfrak m_A B)$.
Enfin, $S/\mathfrak m_R S$ est un anneau local régulier.

\medskip\noindent
Supposons (1), c'est-à-dire que $J$ soit engendré par une suite régulière.
Comme $A = R/I \to B = S/J$ est plat, le
lemme \ref{lemma-perfect-map-ci} s'applique et montre
que $I$ et $J/IS$ sont engendrés par des suites régulières.
On a $\dim(B) = \dim(A) + \dim(B/\mathfrak m_A B)$ et
$\dim(S/IS) = \dim(A) + \dim(S/\mathfrak m_R S)$
(Algèbre, lemme \ref{algebra-lemma-dimension-base-fibre-equals-total}).
Ainsi, $J/IS$ est engendré par
$$
\dim(S/IS) - \dim(S/J) = \dim(S/\mathfrak m_R S) - \dim(B/\mathfrak m_A B)
$$
éléments (Algèbre, lemme \ref{algebra-lemma-one-equation}).
Il s'ensuit que $\Ker(S/\mathfrak m_R S \to B/\mathfrak m_A B)$
est engendré par le même nombre d'éléments (voir ci-dessus).
Par conséquent, $\Ker(S/\mathfrak m_R S \to B/\mathfrak m_A B)$
est engendré par une suite régulière ; voir, par exemple, le
lemme \ref{lemma-quotient-regular-ring}.
On obtient ainsi (2).

\medskip\noindent
Si (2) est vérifiée, alors $I$ et $J/J \cap \mathfrak m_RS$
sont engendrés par des suites régulières. En relevant ces générateurs
(voir ci-dessus), puis en utilisant la platitude de $R/I \to S/IS$
et le lemme de Grothendieck
(Algèbre, lemme \ref{algebra-lemma-grothendieck-regular-sequence}),
on voit que $J/IS$ est engendré par une suite régulière dans $S/IS$.
Le lemme \ref{lemma-perfect-map-ci} montre alors que $J$
est engendré par une suite régulière, d'où (1).
\end{proof}

\begin{definition}
\label{definition-lci}
Soit $A$ un anneau noethérien.
\begin{enumerate}
\item Si $A$ est local, on dit que $A$ est une {\it intersection complète}
si son complété est une intersection complète au sens ci-dessus.
\item En général, on dit que $A$ est une {\it intersection complète locale}
si tous ses anneaux locaux sont des intersections complètes.
\end{enumerate}
\end{definition}

\noindent
Nous vérifierons plus loin que cela ne contredit pas la terminologie
introduite dans
Algèbre, définitions \ref{algebra-definition-lci-field} et
\ref{algebra-definition-lci-local-ring}.
Mais montrons d'abord que cette définition est cohérente : si $A$ est un
anneau local noethérien qui est une intersection complète, alors $A$ est une
intersection complète locale, c'est-à-dire que tous ses anneaux locaux sont
des intersections complètes.

\begin{lemma}
\label{lemma-ci-good}
Soit $(A, \mathfrak m)$ un anneau local noethérien. Soit
$\mathfrak p \subset A$ un idéal premier. Si $A$ est une intersection
complète, alors $A_\mathfrak p$ est aussi une intersection complète.
\end{lemma}

\begin{proof}
Choisissons un idéal premier $\mathfrak q$ de $A^\wedge$ au-dessus de
$\mathfrak p$
(cela est possible puisque $A \to A^\wedge$ est fidèlement plat d'après
Algèbre, lemme \ref{algebra-lemma-completion-faithfully-flat}).
Alors $A_\mathfrak p \to (A^\wedge)_\mathfrak q$ est un homomorphisme
local plat. La proposition \ref{proposition-avramov} montre donc que
$A_\mathfrak p$ est une intersection complète si et seulement si
$(A^\wedge)_\mathfrak q$ est une intersection complète. Il suffit ainsi
de démontrer le lemme lorsque $A$ est complet (c'est l'étape essentielle
de la démonstration).

\medskip\noindent
Supposons $A$ complet. Par définition, on peut écrire
$A = R/(f_1, \ldots, f_r)$ pour une suite régulière
$f_1, \ldots, f_r$ dans un anneau local régulier $R$.
Soit $\mathfrak q \subset R$ l'idéal premier correspondant à
$\mathfrak p$.
On observe que $f_1, \ldots, f_r \in \mathfrak q$ et que
$A_\mathfrak p = R_\mathfrak q/(f_1, \ldots, f_r)R_\mathfrak q$.
Ainsi, $A_\mathfrak p$ est une intersection complète d'après le
lemme \ref{lemma-quotient-regular-ring-by-regular-sequence}.
\end{proof}

\begin{lemma}
\label{lemma-check-lci-at-maximal-ideals}
Soit $A$ un anneau noethérien. Alors $A$ est une intersection complète
locale si et seulement si $A_\mathfrak m$ est une intersection complète
pour tout idéal maximal $\mathfrak m$ de $A$.
\end{lemma}

\begin{proof}
Cela résulte immédiatement du lemme \ref{lemma-ci-good} et des définitions.
\end{proof}

\begin{lemma}
\label{lemma-check-lci-agrees}
Soit $S$ une algèbre de type fini sur un corps $k$.
\begin{enumerate}
\item pour un idéal premier $\mathfrak q \subset S$, l'anneau local
$S_\mathfrak q$
est une intersection complète au sens de la définition
\ref{algebra-definition-lci-local-ring} d'Algèbre,
si et seulement si $S_\mathfrak q$ est une intersection
complète au sens de la définition \ref{definition-lci}, et
\item $S$ est une intersection complète locale au sens de la définition
\ref{algebra-definition-lci-field} d'Algèbre,
si et seulement si $S$ est une intersection complète
locale au sens de la définition \ref{definition-lci}.
\end{enumerate}
\end{lemma}

\begin{proof}
Démontrons (1). Soit $k[x_1, \ldots, x_n] \to S$ une surjection.
Soit $\mathfrak p \subset k[x_1, \ldots, x_n]$ l'idéal premier
correspondant à $\mathfrak q$.
Soit $I \subset k[x_1, \ldots, x_n]$ le noyau de cette surjection.
L'application $k[x_1, \ldots, x_n]_\mathfrak p \to S_\mathfrak q$
est surjective, de noyau $I_\mathfrak p$. Or
$k[x_1, \ldots, x_n]$ est un anneau régulier d'après
Algèbre, proposition
\ref{algebra-proposition-finite-gl-dim-polynomial-ring}.
L'équivalence des deux notions de (1) résulte donc de la combinaison du
lemme \ref{lemma-quotient-regular-ring}
avec Algèbre, lemme \ref{algebra-lemma-lci-local}.

\medskip\noindent
Après avoir démontré (1), l'équivalence de (2) résulte de la
définition et d'Algèbre, lemme \ref{algebra-lemma-lci-global}.
\end{proof}

\begin{lemma}
\label{lemma-avramov}
Soit $A \to B$ un homomorphisme local plat d'anneaux locaux noethériens.
Les propriétés suivantes sont équivalentes :
\begin{enumerate}
\item $B$ est une intersection complète,
\item $A$ et $B/\mathfrak m_A B$ sont des intersections complètes.
\end{enumerate}
\end{lemma}

\begin{proof}
Maintenant que la définition est cohérente, il s'agit d'une reformulation
immédiate de la proposition (non triviale) \ref{proposition-avramov}.
\end{proof}








\section{Homomorphismes locaux d'intersection complète}
\label{section-lci-homomorphisms}

\noindent
Soit $A \to B$ un homomorphisme local d'anneaux locaux noethériens complets.
Une conséquence du théorème de structure de Cohen est que l'on peut trouver un
diagramme commutatif
$$
\xymatrix{
S \ar[r] & B \\
& A \ar[lu] \ar[u]
}
$$
d'anneaux locaux noethériens complets tel que
$S \to B$ soit surjectif, $A \to S$ soit plat et que
$S/\mathfrak m_A S$ soit un anneau local régulier. Cela résulte de
Compléments sur l'algèbre, lemme
\ref{more-algebra-lemma-embed-map-Noetherian-complete-local-rings}.
Disons (provisoirement) que $A \to S \to B$ est une {\it bonne factorisation}
de $A \to B$ si $S$ est un anneau local noethérien,
$A \to S \to B$ sont des homomorphismes locaux d'anneaux,
$S \to B$ est surjectif, $A \to S$ est plat et $S/\mathfrak m_AS$ est régulier.
Disons que $A \to B$ est un
{\it homomorphisme d'intersection complète}
s'il existe une bonne factorisation $A \to S \to B$
telle que le noyau de $S \to B$ soit engendré par une suite régulière.
Le lemme suivant montre que cette notion ne dépend pas
du choix du diagramme.

\begin{lemma}
\label{lemma-ci-map-well-defined}
Soit $A \to B$ un homomorphisme local d'anneaux locaux noethériens complets.
Les assertions suivantes sont équivalentes :
\begin{enumerate}
\item pour une certaine bonne factorisation $A \to S \to B$, le noyau de
$S \to B$ est engendré par une suite régulière, et
\item pour toute bonne factorisation $A \to S \to B$, le noyau de
$S \to B$ est engendré par une suite régulière.
\end{enumerate}
\end{lemma}

\begin{proof}
Soit $A \to S \to B$ une bonne factorisation.
Comme $B$ est complet, on obtient une factorisation
$A \to S^\wedge \to B$ où $S^\wedge$ est le complété de $S$.
Remarquons qu'il s'agit encore d'une bonne factorisation :
l'homomorphisme d'anneaux $S \to S^\wedge$ est plat
(Algèbre, lemme \ref{algebra-lemma-completion-flat}),
donc $A \to S^\wedge$ est plat.
L'anneau $S^\wedge/\mathfrak m_A S^\wedge = (S/\mathfrak m_A S)^\wedge$
est régulier puisque $S/\mathfrak m_A S$ est régulier
(Compléments sur l'algèbre, lemme \ref{more-algebra-lemma-completion-regular}).
Soit $f_1, \ldots, f_r$ une suite minimale de générateurs
de $\Ker(S \to B)$. Nous utiliserons sans autre mention
le fait qu'un idéal d'un anneau local noethérien est engendré par une suite
régulière si et seulement si tout système minimal de générateurs est une
suite régulière. Remarquons que $f_1, \ldots, f_r$
est une suite régulière dans $S$ si et seulement si $f_1, \ldots, f_r$
est une suite régulière dans le complété $S^\wedge$, d'après
Algèbre, lemme \ref{algebra-lemma-flat-increases-depth}.
De plus, on a
$$
S^\wedge/(f_1, \ldots, f_r)R^\wedge =
(S/(f_1, \ldots, f_n))^\wedge = B^\wedge = B
$$
car $B$ est complet pour la topologie $\mathfrak m_B$-adique (la première égalité
résulte d'Algèbre, lemme \ref{algebra-lemma-completion-tensor}).
Ainsi, le noyau de $S \to B$ est engendré par une suite régulière
si et seulement si le noyau de $S^\wedge \to B$ est engendré par une
suite régulière.
Il suffit donc de considérer les bonnes factorisations pour lesquelles $S$ est complet.

\medskip\noindent
Supposons données deux factorisations $A \to S \to B$ et
$A \to S' \to B$, avec $S$ et $S'$ complets. D'après
Compléments sur l'algèbre, lemme \ref{more-algebra-lemma-dominate-two-surjections},
l'anneau $S \times_B S'$ est un anneau local noethérien complet.
Par conséquent, en appliquant Compléments sur l'algèbre, lemme
\ref{more-algebra-lemma-embed-map-Noetherian-complete-local-rings},
on peut choisir une bonne factorisation $A \to S'' \to S \times_B S'$
avec $S''$ complet. Il suffit donc de montrer ceci :
si $A \to S' \to S \to B$ sont de bonnes factorisations comparables,
alors $\Ker(S \to B)$ est engendré par une suite régulière
si et seulement si $\Ker(S' \to B)$ est engendré par une suite régulière.

\medskip\noindent
Soient $A \to S' \to S \to B$ de bonnes factorisations comparables.
Tout d'abord, puisque $S'/\mathfrak m_R S' \to S/\mathfrak m_R S$ est
une surjection d'anneaux locaux réguliers, son noyau est engendré
par une suite régulière
$\overline{x}_1, \ldots, \overline{x}_c \in
\mathfrak m_{S'}/\mathfrak m_R S'$
que l'on peut compléter en un système régulier de paramètres de
l'anneau local régulier $S'/\mathfrak m_R S'$ ; voir
(Algèbre, lemme \ref{algebra-lemma-regular-quotient-regular}).
Posons $I = \Ker(S' \to S)$. Comme $S$ est plat sur $R$, on a
$$
I/\mathfrak m_R I =
\Ker(S'/\mathfrak m_R S' \to S/\mathfrak m_R S) =
(\overline{x}_1, \ldots, \overline{x}_c).
$$
Choisissons des relèvements $x_1, \ldots, x_c \in I$. Ceux-ci forment une suite régulière
qui engendre $I$, puisque $S'$ est plat sur $R$ ; voir
Algèbre, lemme \ref{algebra-lemma-grothendieck-regular-sequence}.

\medskip\noindent
On en déduit que, si $\Ker(S \to B)$ est lui aussi engendré par une
suite régulière, il en va de même de $\Ker(S' \to B)$ ; voir
Compléments sur l'algèbre, lemmes
\ref{more-algebra-lemma-join-koszul-regular-sequences} et
\ref{more-algebra-lemma-noetherian-finite-all-equivalent}.

\medskip\noindent
Réciproquement, supposons que $J = \Ker(S' \to B)$ soit engendré
par une suite régulière. Comme les générateurs $x_1, \ldots, x_c$
de $I$ s'envoient sur des éléments linéairement indépendants de
$\mathfrak m_{S'}/\mathfrak m_{S'}^2$, on voit que
$I/\mathfrak m_{S'}I \to J/\mathfrak m_{S'}J$ est injectif.
Il existe donc un système minimal de générateurs
$x_1, \ldots, x_c, y_1, \ldots, y_d$ de $J$.
Alors $x_1, \ldots, x_c, y_1, \ldots, y_d$ est une suite régulière,
et il s'ensuit que les images de $y_1, \ldots, y_d$ dans $S$
forment une suite régulière qui engendre $\Ker(S \to B)$.
Ceci achève la démonstration du lemme.
\end{proof}

\noindent
Dans la proposition suivante, remarquons que la condition d'annulation des
Tor est notamment satisfaite si $B$ est de dimension de Tor finie sur $A$, et
donc, en particulier, si $B$ est plat sur $A$.

\begin{proposition}
\label{proposition-avramov-map}
Soit $A \to B$ un homomorphisme local d'anneaux locaux noethériens.
Les assertions suivantes sont alors équivalentes :
\begin{enumerate}
\item $B$ est une intersection complète et
$\text{Tor}^A_p(B, A/\mathfrak m_A)$ n'est non nul que pour un nombre fini de $p$,
\item $A$ est une intersection complète et
$A^\wedge \to B^\wedge$ est un homomorphisme d'intersection complète
au sens défini ci-dessus.
\end{enumerate}
\end{proposition}

\begin{proof}
Soit $F_\bullet \to A/\mathfrak m_A$ une résolution par des $A$-modules
libres de type fini. Remarquons que
$\text{Tor}^A_p(B, A/\mathfrak m_A)$
est la $p$-ième homologie du complexe $F_\bullet \otimes_A B$.
Soit $F_\bullet^\wedge = F_\bullet \otimes_A A^\wedge$ le complété.
Alors $F_\bullet^\wedge$ est une résolution de $A^\wedge/\mathfrak m_{A^\wedge}$
par des $A^\wedge$-modules libres de type fini (car $A \to A^\wedge$ est plat et le complété
est exact sur les modules finis ; voir
Algèbre, lemmes \ref{algebra-lemma-completion-tensor} et
\ref{algebra-lemma-completion-flat}).
Il s'ensuit que
$$
F_\bullet^\wedge \otimes_{A^\wedge} B^\wedge =
F_\bullet \otimes_A B \otimes_B B^\wedge
$$
La platitude de $B \to B^\wedge$ donne alors
$$
\text{Tor}^{A^\wedge}_p(B^\wedge, A^\wedge/\mathfrak m_{A^\wedge}) =
\text{Tor}^A_p(B, A/\mathfrak m_A) \otimes_B B^\wedge
$$
On voit ainsi que la condition de (1) portant sur l'homomorphisme local $A \to B$
est équivalente à la même condition pour l'homomorphisme local
$A^\wedge \to B^\wedge$.
On peut donc supposer que $A$ et $B$ sont des anneaux locaux noethériens complets
(puisque, de toute façon, les autres conditions sont formulées en termes des complétés).

\medskip\noindent
Supposons que $A$ et $B$ soient des anneaux locaux noethériens complets.
Choisissons un diagramme
$$
\xymatrix{
S \ar[r] & B \\
R \ar[u] \ar[r] & A \ar[u]
}
$$
comme dans Compléments sur l'algèbre, lemme
\ref{more-algebra-lemma-embed-map-Noetherian-complete-local-rings}.
Posons $I = \Ker(R \to A)$ et $J = \Ker(S \to B)$.
La proposition résulte alors du lemme \ref{lemma-perfect-map-ci}.
\end{proof}

\begin{remark}
\label{remark-no-good-ci-map}
Il semble difficile de définir une bonne notion d'« homomorphismes locaux
d'intersection complète » pour les homomorphismes entre anneaux noethériens généraux.
La raison en est que, pour un anneau local noethérien $A$, les fibres de
$A \to A^\wedge$ ne sont pas des anneaux d'intersection complète locale.
Ainsi, si $A \to B$ est un homomorphisme local d'anneaux locaux noethériens
et si l'homomorphisme entre les complétés $A^\wedge \to B^\wedge$ est un
homomorphisme d'intersection complète au sens défini ci-dessus,
alors $(A_\mathfrak p)^\wedge \to (B_\mathfrak q)^\wedge$ n'est en général
{\bf pas} un homomorphisme d'intersection complète au sens
défini ci-dessus. Une solution consiste à travailler exclusivement avec des
anneaux noethériens excellents. Plus généralement, on pourrait travailler avec
les anneaux noethériens dont les fibres formelles sont des
intersections complètes ; voir \cite{Rodicio-ci}.
Nous développerons cette théorie dans
Complexes dualisants, section \ref{dualizing-section-formal-fibres}.
\end{remark}

\noindent
Pour terminer cette section, comparons la notion définie ci-dessus
à celle introduite dans
Compléments sur l'algèbre, section \ref{section-lci}.

\begin{lemma}
\label{lemma-well-defined-if-you-can-find-good-factorization}
Considérons un diagramme commutatif
$$
\xymatrix{
S \ar[r] & B \\
& A \ar[lu] \ar[u]
}
$$
d'anneaux locaux noethériens tel que $S \to B$ soit surjectif, $A \to S$ soit plat et
que $S/\mathfrak m_A S$ soit un anneau local régulier. Les assertions suivantes sont équivalentes :
\begin{enumerate}
\item $\Ker(S \to B)$ est engendré par une suite régulière, et
\item $A^\wedge \to B^\wedge$ est un homomorphisme d'intersection complète
au sens défini ci-dessus.
\end{enumerate}
\end{lemma}

\begin{proof}
Démonstration omise. Indication : la démonstration est identique à l'argument donné dans
le premier paragraphe de la démonstration du lemme \ref{lemma-ci-map-well-defined}.
\end{proof}

\begin{lemma}
\label{lemma-finite-type-lci-map}
Soit $A$ un anneau noethérien.
Soit $A \to B$ un homomorphisme d'anneaux de type fini.
Les assertions suivantes sont équivalentes :
\begin{enumerate}
\item $A \to B$ est une intersection complète locale au sens de
Compléments sur l'algèbre, définition
\ref{more-algebra-definition-local-complete-intersection},
\item pour tout idéal premier $\mathfrak q \subset B$, en posant
$\mathfrak p = A \cap \mathfrak q$, l'homomorphisme d'anneaux
$(A_\mathfrak p)^\wedge \to (B_\mathfrak q)^\wedge$ est
un homomorphisme d'intersection complète au sens défini ci-dessus.
\end{enumerate}
\end{lemma}

\begin{proof}
Choisissons une surjection $R = A[x_1, \ldots, x_n] \to B$.
Remarquons que $A \to R$ est plat à fibres régulières.
Soit $I$ le noyau de $R \to B$.
Supposons (2). On voit alors que
$I$ est localement engendré par une suite régulière,
d'après
le lemme \ref{lemma-well-defined-if-you-can-find-good-factorization}
et
Algèbre, lemme \ref{algebra-lemma-regular-sequence-in-neighbourhood}.
Autrement dit, (1) est vérifiée.
Réciproquement, supposons (1). Après localisation sur $R$ et $B$,
on peut alors supposer que $I$ est engendré par une suite régulière au sens de Koszul.
D'après Compléments sur l'algèbre, lemme
\ref{more-algebra-lemma-noetherian-finite-all-equivalent},
on trouve que $I$ est localement engendré par une suite régulière.
Ainsi, (2) est vérifiée d'après
le lemme \ref{lemma-well-defined-if-you-can-find-good-factorization}.
Certains détails sont omis.
\end{proof}

\begin{lemma}
\label{lemma-avramov-map-finite-type}
Soit $A$ un anneau noethérien. Soit $A \to B$ un homomorphisme d'anneaux de type fini
tel que l'image de $\Spec(B) \to \Spec(A)$ contienne tous les points fermés
de $\Spec(A)$. Les assertions suivantes sont alors équivalentes :
\begin{enumerate}
\item $B$ est une intersection complète et $A \to B$ est de dimension de
Tor finie,
\item $A$ est une intersection complète et $A \to B$ est une intersection complète
locale au sens de Compléments sur l'algèbre, définition
\ref{more-algebra-definition-local-complete-intersection}.
\end{enumerate}
\end{lemma}

\begin{proof}
Il s'agit d'une reformulation de la proposition \ref{proposition-avramov-map}
au moyen du lemme \ref{lemma-finite-type-lci-map}.
Nous omettons les détails.
\end{proof}








\section{Homomorphismes d'anneaux lisses et diagonales}
\label{section-smooth-diagonal-perfect}

\noindent
Dans cette section, nous utilisons les résultats précédents pour caractériser les homomorphismes d'anneaux lisses comme
ceux dont la diagonale est parfaite.

\begin{lemma}
\label{lemma-local-perfect-diagonal}
Soit $A \to B$ un homomorphisme local d'anneaux locaux noethériens tel que
$B$ soit plat et essentiellement de type fini sur $A$. Si
$$
B \otimes_A B \longrightarrow B
$$
est un homomorphisme d'anneaux parfait, c'est-à-dire si $B$ est de dimension de Tor finie sur
$B \otimes_A B$, alors $B$ est la localisation d'une $A$-algèbre lisse.
\end{lemma}

\begin{proof}
Comme $B$ est essentiellement de type fini sur $A$, il en va de même de $B \otimes_A B$, et
en particulier $B \otimes_A B$ est noethérien. Par conséquent, le quotient $B$ de
$B \otimes_A B$ est pseudo-cohérent sur $B \otimes_A B$
(Compléments sur l'algèbre, lemme \ref{more-algebra-lemma-Noetherian-pseudo-coherent}),
ce qui explique pourquoi la perfection de l'homomorphisme d'anneaux (Compléments sur l'algèbre, définition
\ref{more-algebra-definition-pseudo-coherent-perfect}) équivaut à la
condition de dimension de Tor finie.

\medskip\noindent
On peut écrire $B = R/K$, où $R$ est la localisation de $A[x_1, \ldots, x_n]$
en un idéal premier et où $K \subset R$ est un idéal. Notons
$\mathfrak m \subset R \otimes_A R$ l'idéal maximal qui est l'image réciproque
de l'idéal maximal de $B$ par la surjection
$R \otimes_A R \to B \otimes_A B \to B$. On a alors des surjections
$$
(R \otimes_A R)_\mathfrak m \to (B \otimes_A B)_\mathfrak m \to B
$$
et donc des idéaux $I \subset J \subset (R \otimes_A R)_\mathfrak m$
comme dans le lemme \ref{lemma-injective}. On en déduit que
$I/\mathfrak m I \to J/\mathfrak m J$ est injectif.

\medskip\noindent
Écrivons $K = (f_1, \ldots, f_r)$ avec $r$ minimal. On peut supposer, et on le fera,
que $f_i \in R$ est l'image d'un élément de $A[x_1, \ldots, x_n]$ que l'on
note également $f_i$. Remarquons que $I$ est engendré
par $f_1 \otimes 1, \ldots, f_r \otimes 1$ et
$1 \otimes f_1, \ldots, 1 \otimes f_r$. Nous affirmons que c'est un système minimal
de générateurs de $I$. En effet, si $\kappa$ est le corps résiduel commun
de $R$, $B$, $(R \otimes_A R)_\mathfrak m$ et $(B \otimes_A B)_\mathfrak m$,
alors on dispose d'un homomorphisme
$R \otimes_A R \to R \otimes_A \kappa \oplus \kappa \otimes_A R$
qui se factorise par $(R \otimes_A R)_\mathfrak m$. Comme $B$ est
plat sur $A$ et que l'on a la suite exacte courte
$0 \to K \to R \to B \to 0$, on voit que
$K \otimes_A \kappa \subset R \otimes_A \kappa$ ; voir
Algèbre, lemme \ref{algebra-lemma-flat-tor-zero}.
En restreignant ainsi l'homomorphisme
$(R \otimes_A R)_\mathfrak m \to R \otimes_A \kappa \oplus \kappa \otimes_A R$
à $I$, on obtient un homomorphisme
$$
I \to K \otimes_A \kappa \oplus \kappa \otimes_A K \to
K \otimes_B \kappa \oplus \kappa \otimes_B K.
$$
Les éléments
$f_1 \otimes 1, \ldots, f_r \otimes 1, 1 \otimes f_1, \ldots, 1 \otimes f_r$
s'envoient sur une base du but de cet homomorphisme, puisque, d'après le lemme de Nakayama
(Algèbre, lemme \ref{algebra-lemma-NAK}),
$f_1, \ldots, f_r$ s'envoient sur une base de $K \otimes_B \kappa$.
Ceci démontre notre assertion.

\medskip\noindent
L'idéal $J$ est engendré par $f_1 \otimes 1, \ldots, f_r \otimes 1$
et par les éléments $x_1 \otimes 1 - 1 \otimes x_1, \ldots,
x_n \otimes 1 - 1 \otimes x_n$ (pour la démonstration, il suffit de
voir que ces éléments appartiennent à l'idéal $J$).
On peut alors écrire
$$
f_i \otimes 1 - 1 \otimes f_i =
\sum g_{ij} (x_j \otimes 1 - 1 \otimes x_j)
$$
pour certains $g_{ij}$ dans $(R \otimes_A R)_\mathfrak m$. C'est un fait général
concernant les éléments de $A[x_1, \ldots, x_n]$, dont nous omettons la démonstration.
Notons $a_{ij} \in \kappa$ l'image de $g_{ij}$. Un autre calcul
montre que $a_{ij}$ est l'image de $\partial f_i / \partial x_j$ dans $\kappa$.
L'injectivité de $I/\mathfrak m I \to J/\mathfrak m J$ et les remarques
précédentes imposent à la matrice $(a_{ij})$ d'être de rang maximal $r$.
Posons
$$
C = A[x_1, \ldots, x_n]/(f_1, \ldots, f_r)
$$
et considérons le complexe cotangent naïf
$\NL_{C/A} \cong (C^{\oplus r} \to C^{\oplus n})$
où l'homomorphisme est donné par la matrice des dérivées partielles.
Ainsi, $\NL_{C/A} \otimes_A B$
est isomorphe à un $B$-module libre de rang $n - r$ placé en degré $0$.
Ainsi, $C_g$ est lisse sur $A$ pour un certain $g \in C$ dont l'image dans $B$ est une unité ; voir Algèbre, lemme \ref{algebra-lemma-smooth-at-point}.
Ceci achève la démonstration.
\end{proof}

\begin{lemma}
\label{lemma-perfect-diagonal}
Soit $A \to B$ un homomorphisme d'anneaux noethériens, plat et de type fini. Si
$$
B \otimes_A B \longrightarrow B
$$
est un homomorphisme d'anneaux parfait, c'est-à-dire si $B$ est de dimension de Tor finie sur
$B \otimes_A B$, alors $B$ est une $A$-algèbre lisse.
\end{lemma}

\begin{proof}
Cela résulte du lemme \ref{lemma-local-perfect-diagonal}
et des propriétés générales des homomorphismes d'anneaux lisses ; voir
Algèbre, lemmes \ref{algebra-lemma-smooth-at-point} et
\ref{algebra-lemma-locally-smooth}.
On peut aussi modifier légèrement la démonstration du
lemme \ref{lemma-local-perfect-diagonal} pour démontrer
le présent lemme.
\end{proof}






\section{Liberté du module conormal}
\label{section-freeness-conormal}

\noindent
Les résolutions de Tate et les dérivations sur celles-ci permettent de démontrer
des versions (plus fortes) des résultats de cette section ; voir \cite{Iyengar}.
Deux références plus élémentaires sont
\cite{Vasconcelos} et \cite{Ferrand-lci}.

\begin{lemma}
\label{lemma-free-summand-in-ideal-finite-proj-dim}
\begin{reference}
\cite{Vasconcelos}
\end{reference}
Soit $R$ un anneau local noethérien. Soit $I \subset R$ un idéal
de dimension projective finie sur $R$. Si $F \subset I/I^2$ est un
facteur direct isomorphe à $R/I$, alors il existe un élément non diviseur de zéro
$x \in I$ tel que l'image de $x$ dans $I/I^2$ engendre $F$.
\end{lemma}

\begin{proof}
Par hypothèse, on peut choisir une résolution libre finie
$$
0 \to R^{\oplus n_e} \to R^{\oplus n_{e-1}} \to \ldots \to
R^{\oplus n_1} \to R \to R/I \to 0
$$
Alors $\varphi_1 : R^{\oplus n_1} \to R$ est de rang $1$, et
on voit que $I$ contient un élément non diviseur de zéro $y$, d'après
Algèbre, proposition \ref{algebra-proposition-what-exact}.
Soient $\mathfrak p_1, \ldots, \mathfrak p_n$ les idéaux premiers associés
à $R$ ; voir Algèbre, lemme \ref{algebra-lemma-finite-ass}.
Soit $I^2 \subset J \subset I$ un idéal tel que $J/I^2 = F$.
Alors $J \not \subset \mathfrak p_i$ pour tout $i$,
car $y^2 \in J$ et $y^2 \not \in \mathfrak p_i$ ; voir
Algèbre, lemme \ref{algebra-lemma-ass-zero-divisors}.
D'après le lemme de Nakayama (Algèbre, lemme \ref{algebra-lemma-NAK}),
on a $J \not \subset \mathfrak m J + I^2$.
D'après Algèbre, lemme \ref{algebra-lemma-silly},
on peut choisir $x \in J$, $x \not \in \mathfrak m J + I^2$ et
$x \not \in \mathfrak p_i$ pour $i = 1, \ldots, n$.
Alors $x$ est un élément non diviseur de zéro, et l'image
de $x$ dans $I/I^2$ engendre (d'après le lemme de Nakayama)
le facteur direct $J/I^2 \cong R/I$.
\end{proof}

\begin{lemma}
\label{lemma-vasconcelos}
\begin{reference}
Version locale de \cite[Theorem 1.1]{Vasconcelos}
\end{reference}
Soit $R$ un anneau local noethérien. Soit $I \subset R$ un idéal
de dimension projective finie sur $R$. Si $F \subset I/I^2$
est un facteur direct libre de rang $r$, alors il existe une suite régulière
$x_1, \ldots, x_r \in I$ telle que $x_1 \bmod I^2, \ldots, x_r \bmod I^2$
engendrent $F$.
\end{lemma}

\begin{proof}
Si $r = 0$, il n'y a rien à démontrer. Supposons $r > 0$. On peut choisir
$x \in I$ tel que $x$ soit un élément non diviseur de zéro et que $x \bmod I^2$
engendre un facteur direct de $F$ isomorphe à $R/I$ ; voir
le lemme \ref{lemma-free-summand-in-ideal-finite-proj-dim}.
Considérons l'anneau $R' = R/(x)$ et l'idéal $I' = I/(x)$.
Bien entendu, $R'/I' = R/I$. La suite exacte courte
$$
0 \to R/I \xrightarrow{x} I/xI \to I' \to 0
$$
est scindée, car l'homomorphisme $I/xI \to I/I^2$ envoie $xR/xI$
sur un facteur direct. Or $I/xI = I \otimes_R^\mathbf{L} R'$ est de
dimension projective finie sur $R'$ ; voir
Compléments sur l'algèbre, lemmes \ref{more-algebra-lemma-perfect-module} et
\ref{more-algebra-lemma-pull-perfect}.
Par conséquent, le facteur direct $I'$ est de dimension projective finie sur $R'$.
D'autre part, on a la suite exacte courte
$0 \to xR/xI \to I/I^2 \to I'/(I')^2 \to 0$, et on en déduit que
$I'/(I')^2$ possède comme facteur direct libre $F' = F/(R/I \cdot x)$
de rang $r - 1$. Par récurrence sur $r$,
on peut choisir une suite régulière $x'_2, \ldots, x'_r \in I'$
dont les classes engendrent librement $F'$.
Si $x_1 = x$ et si $x_2, \ldots, x_r$ sont des relèvements quelconques de
$x'_1, \ldots, x'_r$ dans $R$, on voit alors que le lemme est vérifié.
\end{proof}

\begin{proposition}
\label{proposition-regular-ideal}
\begin{reference}
Variante de \cite[Corollary 1]{Vasconcelos}. Voir aussi
\cite{Iyengar} et \cite{Ferrand-lci}.
\end{reference}
Soit $R$ un anneau noethérien. Soit $I \subset R$ un idéal
de dimension projective finie tel que $I/I^2$ soit
localement libre de type fini sur $R/I$. Alors $I$ est un idéal régulier
(Compléments sur l'algèbre, définition \ref{more-algebra-definition-regular-ideal}).
\end{proposition}

\begin{proof}
D'après Algèbre, lemme \ref{algebra-lemma-regular-sequence-in-neighbourhood},
il suffit de montrer que $I_\mathfrak p \subset R_\mathfrak p$ est engendré
par une suite régulière pour tout $\mathfrak p \supset I$. On peut donc
supposer que $R$ est local. Si $I/I^2$ est de rang $r$, alors le
lemme \ref{lemma-vasconcelos} fournit une suite régulière
$x_1, \ldots, x_r \in I$ qui engendre $I/I^2$. D'après
le lemme de Nakayama (Algèbre, lemme \ref{algebra-lemma-NAK}),
on en déduit que $I$ est engendré par $x_1, \ldots, x_r$.
\end{proof}

\noindent
Pour tout homomorphisme d'anneaux d'intersection complète locale $A \to B$,
le complexe cotangent naïf $\NL_{B/A}$ est parfait
d'amplitude de Tor dans $[-1, 0]$ ; voir
Compléments sur l'algèbre, lemme \ref{more-algebra-lemma-lci-NL}.
À l'aide de ce qui précède, on peut montrer que cela caractérise parfois
les homomorphismes d'intersection complète locale.

\begin{lemma}
\label{lemma-perfect-NL-lci}
Soit $A \to B$ un homomorphisme d'anneaux parfait (Compléments sur l'algèbre, définition
\ref{more-algebra-definition-pseudo-coherent-perfect})
entre anneaux noethériens. Les assertions suivantes sont alors équivalentes :
\begin{enumerate}
\item $\NL_{B/A}$ est d'amplitude de Tor dans $[-1, 0]$,
\item $\NL_{B/A}$ est un objet parfait de $D(B)$
d'amplitude de Tor dans $[-1, 0]$, et
\item $A \to B$ est une intersection complète locale
(Compléments sur l'algèbre, définition
\ref{more-algebra-definition-local-complete-intersection}).
\end{enumerate}
\end{lemma}

\begin{proof}
Écrivons $B = A[x_1, \ldots, x_n]/I$. Alors $\NL_{B/A}$ est représenté par
le complexe
$$
I/I^2 \longrightarrow \bigoplus B \text{d}x_i
$$
de $B$-modules dans lequel $I/I^2$ est placé en degré $-1$. Puisque le terme de
degré $0$ est libre de type fini, ce complexe est d'amplitude de Tor dans $[-1, 0]$ si et
seulement si $I/I^2$ est un $B$-module plat ; voir
Compléments sur l'algèbre, lemme \ref{more-algebra-lemma-last-one-flat}.
Comme $I/I^2$ est un $B$-module fini et que $B$ est noethérien, cela équivaut
à ce que $I/I^2$ soit un $B$-module localement libre de type fini
(Algèbre, lemme \ref{algebra-lemma-finite-projective}).
L'équivalence de (1) et (2) est donc claire. En outre, l'équivalence
de (1) et (3) résulte également de la
proposition \ref{proposition-regular-ideal}
(et du fait qu'un idéal régulier est à la fois un idéal régulier au sens de Koszul
et un idéal quasi-régulier ; voir
Compléments sur l'algèbre, section \ref{more-algebra-section-ideals}).
\end{proof}

\begin{lemma}
\label{lemma-flat-fp-NL-lci}
Soit $A \to B$ un homomorphisme d'anneaux plat et de présentation finie.
Les assertions suivantes sont alors équivalentes :
\begin{enumerate}
\item $\NL_{B/A}$ est d'amplitude de Tor dans $[-1, 0]$,
\item $\NL_{B/A}$ est un objet parfait de $D(B)$
d'amplitude de Tor dans $[-1, 0]$,
\item $A \to B$ est syntomique
(Algèbre, définition \ref{algebra-definition-lci}), et
\item $A \to B$ est une intersection complète locale
(Compléments sur l'algèbre, définition
\ref{more-algebra-definition-local-complete-intersection}).
\end{enumerate}
\end{lemma}

\begin{proof}
L'équivalence de (3) et (4) est Compléments sur l'algèbre, lemme
\ref{more-algebra-lemma-syntomic-lci}.

\medskip\noindent
Si $A \to B$ est syntomique, on peut trouver un diagramme cocartésien
$$
\xymatrix{
B_0 \ar[r] & B \\
A_0 \ar[r] \ar[u] & A \ar[u]
}
$$
tel que $A_0 \to B_0$ soit syntomique et que $A_0$ soit noethérien ; voir
Algèbre, lemmes \ref{algebra-lemma-limit-module-finite-presentation} et
\ref{algebra-lemma-colimit-lci}. D'après le lemme \ref{lemma-perfect-NL-lci},
on voit que $\NL_{B_0/A_0}$ est parfait d'amplitude de Tor dans $[-1, 0]$.
D'après Compléments sur l'algèbre, lemme \ref{more-algebra-lemma-base-change-NL-flat},
on en déduit qu'il en va de même de
$\NL_{B/A} = \NL_{B_0/A_0} \otimes_{B_0}^\mathbf{L} B$ (voir
aussi Compléments sur l'algèbre, lemmes \ref{more-algebra-lemma-pull-tor-amplitude} et
\ref{more-algebra-lemma-pull-perfect}).
Ceci montre que (3) implique (2).

\medskip\noindent
Supposons (1). D'après Compléments sur l'algèbre, lemme
\ref{more-algebra-lemma-base-change-NL-flat},
pour tout homomorphisme d'anneaux $A \to k$, où
$k$ est un corps, on voit que $\NL_{B \otimes_A k/k}$ est
d'amplitude de Tor dans $[-1, 0]$ (voir
Compléments sur l'algèbre, lemme \ref{more-algebra-lemma-pull-tor-amplitude}).
Le lemme \ref{lemma-perfect-NL-lci} montre donc que $k \to B \otimes_A k$ est
un homomorphisme d'intersection complète locale. Ainsi, $A \to B$
est syntomique par définition. Ceci montre que (1) implique (3)
et achève la démonstration.
\end{proof}


\section{Complexes de Koszul et résolutions de Tate}

\label{section-koszul-vs-tate}

\noindent
Dans cette section, nous « relevons » le résultat de Compléments sur l'algèbre, lemme \ref{more-algebra-lemma-sequence-Koszul-complexes}, à la catégorie des algèbres différentielles graduées munies de puissances divisées compatibles avec la structure différentielle graduée (attention : nous représentons ici les complexes de Koszul comme des complexes de chaînes, tandis que dans la référence citée ils sont considérés comme des complexes de cochaînes).

\medskip\noindent
Soit $R$ un anneau. Soit $I \subset R$ l'idéal engendré par $f_1, \ldots, f_r \in R$. Pour $n \geq 1$, nous notons
$$
K_n = K_{n, \bullet} = R\langle \xi_1, \ldots, \xi_r\rangle
$$
l'algèbre différentielle graduée de Koszul, où $\xi_i$ est de degré $1$ et $\text{d}(\xi_i) = f_i^n$. Elle porte une unique structure de puissances divisées (au sens de la définition \ref{definition-divided-powers-dga}) ; voir l'exemple \ref{example-adjoining-odd}. Pour $m > n$, l'application de transition $K_m \to K_n$ est l'homomorphisme d'algèbres différentielles graduées compatible avec les puissances divisées qui envoie $\xi_i$ sur $f_i^{m - n}\xi_i$.

\begin{lemma}

\label{lemma-lift-tate-to-koszul}

Dans la situation ci-dessus, si $R$ est noethérien, alors, pour tout $n$, il existe un entier $N \geq n$ et des applications
$$
K_N \to A \to R/(f_1^N, \ldots, f_r^N)\quad\text{et}\quad A \to K_n
$$
avec les propriétés suivantes

\begin{enumerate}

\item  $(A, \text{d}, \gamma)$ est comme dans la définition \ref{definition-divided-powers-dga},
\item  $A \to R/(f_1^N, \ldots, f_r^N)$ est un quasi-isomorphisme,
\item la composition $K_N \to A \to R/(f_1^N, \ldots, f_r^N)$
est l'application canonique,
\item la composition $K_N \to A \to K_n$ est l'application de transition,
\item  $A_0 = R \to R/(f_1^N, \ldots, f_r^N)$ est la surjection canonique,
\item  $A$ est une algèbre polynomiale graduée à puissances divisées sur $R$
ayant un nombre fini de générateurs dans chaque degré ;
\item  $A \to K_n$ est un homomorphisme d'algèbres différentielles graduées sur $R$, compatible avec les puissances divisées, qui induit l'application canonique
$R/(f_1^N, \ldots, f_r^N) \to R/(f_1^n, \ldots, f_r^n)$ en homologie de degré $0$.

\end{enumerate}

La condition (6) signifie que $A$ est construit à partir de $A_0$ en adjoignant successivement une famille finie de variables $T$ dans chaque degré
$> 0$, comme dans l'exemple \ref{example-adjoining-odd} ou \ref{example-adjoining-even}.

\end{lemma}

\begin{proof}
Fixons $n$. Si $r = 0$, nous pouvons simplement prendre $A = R$. Supposons $r > 0$. D'après Compléments sur l'algèbre, lemme \ref{more-algebra-lemma-sequence-Koszul-complexes} (traduit dans le langage des complexes de chaînes), nous pouvons choisir
$$
n_{r} > n_{r - 1} > \ldots > n_1 > n_0 = n
$$
de telle sorte que les applications de transition $K_{n_{i + 1}} \to K_{n_i}$ entre algèbres de Koszul (définies ci-dessus) induisent l'application nulle en homologie dans les degrés $> 0$. Nous démontrerons le lemme en prenant $N = n_r$.

\medskip\noindent

Nous construisons $A$ exactement comme dans l'énoncé et la démonstration du lemme \ref{lemma-tate-resolution}. Nous aurons donc
$$
A = \colim A(m),\quad\text{et}\quad
A(0) \to A(1) \to A(2) \to \ldots \to R/(f_1^N, \ldots, f_r^N)
$$
On obtient immédiatement les propriétés (1), (2), (5) et (6). Pour achever la démonstration, nous construirons les homomorphismes de $R$-algèbres
$K_N \to A \to K_n$. Pour ce faire, nous allons construire

\begin{enumerate}

\item un isomorphisme $A(1) \to K_N = K_{n_r}$,
\item une application $A(2) \to K_{n_{r - 1}}$,
\item  $\ldots$

\item une application $A(r) \to K_{n_1}$,
\item une application $A(r + 1) \to K_{n_0} = K_n$,
\item une application $A \to K_n$.

\end{enumerate}

À chacune de ces étapes, l'application construite sera un homomorphisme entre algèbres différentielles graduées munies de puissances divisées compatibles ; elle sera compatible avec la précédente par l'intermédiaire des applications de transition entre algèbres de Koszul et induira l'application canonique évidente en homologie de degré $0$.

\medskip\noindent
Rappelons que $A(0) = R$. Pour $m = 1$, la démonstration du lemme \ref{lemma-tate-resolution} prend $A(1) = R\langle T_1, \ldots, T_r\rangle$, où $T_i$ est de degré $1$ et $\text{d}(T_i) = f_i^N$. En effet, les $f_i^N$ engendrent le noyau de $A(0) \to R/(f_1^N, \ldots, f_r^N)$. Pour définir $A(1) \to K_N = K_{n_r}$, nous utilisons donc l'application
$$
\varphi_1 : A(1) \longrightarrow K_{n_r},\quad T_i \longmapsto \xi_i
$$
qui est un isomorphisme.

\medskip\noindent

Pour $m = 2$, la construction dans la démonstration du lemme \ref{lemma-tate-resolution}
choisit des générateurs $e_1, \ldots, e_t \in \Ker(\text{d} : A(1)_1 \to A(1)_0)$. Elle prend
$A(2) = A(1)\langle T_1, \ldots, T_t\rangle$
comme algèbre polynomiale à puissances divisées, avec $T_i$ de degré $2$
et $\text{d}(T_i) = e_i$. Puisque $\varphi_1(e_i)$ est un cocycle de $K_{n_r}$,
son image dans $K_{n_{r - 1}}$ est un bord, par le choix de $n_r$ et $n_{r - 1}$ effectué ci-dessus. Nous pouvons donc construire le diagramme commutatif
$$
\xymatrix{
A(1) \ar[d] \ar[r]_{\varphi_1} & K_{n_r} \ar[d] \\
A(2) \ar[r]^{\varphi_2} & K_{n_{r - 1}}
}
$$
en envoyant $T_i$ sur un élément de degré $2$ dont le bord est l'image de $\varphi_1(e_i)$. L'application $\varphi_2$ existe et est compatible avec la différentielle et les puissances divisées grâce à la propriété universelle de l'algèbre polynomiale à puissances divisées.

\medskip\noindent
L'algèbre $A(m)$ et l'application $\varphi_m : A(m) \to K_{n_{r + 1 - m}}$ se construisent de la même manière pour $m = 2, \ldots, r$.

\medskip\noindent
Étant donnée l'application $A(r) \to K_{n_1}$, la composée $H_r(A(r)) \to H_r(K_{n_1}) \to H_r(K_{n_0}) \subset (K_{n_0})_r$ est nulle. On peut donc la prolonger en une application $A(r + 1) \to K_{n_0} = K_n$ en envoyant sur zéro les nouveaux générateurs polynomiaux de $A(r + 1)$.

\medskip\noindent

Une fois construite l'application $A(r + 1) \to K_{n_0} = K_n$, on la prolonge de l'unique manière possible à $A(r + 2), A(r + 3), \ldots$, en envoyant les nouveaux générateurs polynomiaux sur zéro. Ceci achève la démonstration.

\end{proof}

\begin{remark}

\label{remark-pro-system-koszul}
Dans la situation ci-dessus, si $R$ est noethérien, nous pouvons choisir par récurrence une suite d'entiers $1 = n_0 < n_1 < n_2 < \ldots $ telle que, pour $i = 1, 2, 3, \ldots$, il existe des applications $K_{n_i} \to A_i \to R/(f_1^{n_i}, \ldots, f_r^{n_i})$ et $A_i \to K_{n_{i - 1}}$ comme dans le lemme \ref{lemma-lift-tate-to-koszul}. Notons $A_{i + 1} \to A_i$ la composée $A_{i + 1} \to K_{n_i} \to A_i$. Le diagramme
$$
\xymatrix{
K_{n_1} \ar[d] &
K_{n_2} \ar[d] \ar[l] &
K_{n_3} \ar[d] \ar[l] &
\ldots \ar[l] \\
A_1 \ar[d] &
A_2 \ar[l] \ar[d] &
A_3 \ar[l] \ar[d] &
\ldots \ar[l] \\
K_1 &
K_{n_1} \ar[l] &
K_{n_2} \ar[l] &
\ldots \ar[l]
}
$$
est alors commutatif. On voit ainsi que les systèmes projectifs $(K_n)$ et $(A_n)$ sont pro-isomorphes dans la catégorie des algèbres différentielles graduées sur $R$ munies de puissances divisées compatibles.
\end{remark}

\begin{lemma}

\label{lemma-compute-cohomology-adjoin-variable}

Soient $(A, \text{d}, \gamma)$, $d \geq 1$, $f \in A_{d - 1}$ et $A\langle T \rangle$ comme dans le lemme \ref{lemma-extend-differential}.

\begin{enumerate}

\item Si $d = 1$, il existe une suite exacte longue
$$
\ldots \to H_0(A) \xrightarrow{f} H_0(A) \to H_0(A\langle T \rangle) \to 0
$$

\item Pour $d = 2$, il existe une suite spectrale bornée
$(E_1)_{i, j} = H_{j - i}(A) \cdot T^{[i]}$
qui converge vers $H_{i + j}(A\langle T \rangle)$. Le différentiel
$(d_1)_{i, j} : H_{j - i}(A) \cdot T^{[i]} \to
H_{j - i + 1}(A) \cdot T^{[i - 1]}$
envoie $\xi \cdot T^{[i]}$ sur la classe de $f \xi \cdot T^{[i - 1]}$.
\item On pourra compléter cet énoncé pour les autres degrés, si nécessaire.

\end{enumerate}

\end{lemma}

\begin{proof}
Pour $d = 1$, nous avons une suite exacte courte de complexes
$$
0 \to A \to A\langle T \rangle \to A \cdot T \to 0
$$
et (1) en résulte immédiatement. Pour $d = 2$, nous considérons $A\langle T \rangle$ comme un complexe de chaînes filtré par les sous-complexes
$$
F^pA\langle T \rangle = \bigoplus\nolimits_{i \leq p} A \cdot T^{[i]}
$$
En appliquant la suite spectrale de Homologie, section \ref{homology-section-filtered-complex} (traduite en complexes de chaînes), on obtient (2).
\end{proof}

\noindent
Le lemme suivant sera nécessaire plus tard.

\begin{lemma}

\label{lemma-construct-some-maps}

Dans la situation ci-dessus, pour tous $n \geq t \geq 1$ il existe un $N > n$
et une application
$$
K_t \longrightarrow K_n \otimes_R K_t
$$
dans la catégorie dérivée des $K_N$-modules différentiels gradués à gauche, dont la composée avec l'application de multiplication est l'application de transition (dans l'un ou l'autre sens).

\end{lemma}

\begin{proof}

Commençons par le cas $r = 1$. Posons $f = f_1$. Écrivons $K_t = R\langle x \rangle$,
$K_n = R\langle y \rangle$, $K_N = R\langle z \rangle$
où $x$, $y$ et $z$ sont de degré $1$, et
$\text{d}(x) = f^t$, $\text{d}(y) = f^n$,
$\text{d}(z) = f^N$. Pour tout $N > t$, nous affirmons qu'il existe un quasi-isomorphisme
$$
B_{N, t} = R\langle x, z, u \rangle
\longrightarrow
K_t,\quad
x \mapsto x,\quad
z \mapsto f^{N - t}x,\quad
u \mapsto 0
$$
Ici, le membre de gauche désigne l'algèbre polynomiale à puissances divisées aux variables $x$ et $z$ de degré $1$ et $u$ de degré $2$, munie de
$\text{d}(x) = f^t$, $\text{d}(z) = f^N$,
$\text{d}(u) = z - f^{N - t}x$. Pour démontrer l'assertion, remarquons que les trois sous-modules suivants de
$H_*(R\langle x, z\rangle)$ sont les mêmes

\begin{enumerate}

\item le noyau de $H_*(R\langle x, z\rangle) \to H_*(K_t)$,
\item l'image de
$z - f^{N - t}x : H_*(R\langle x, z\rangle) \to H_*(R\langle x, z\rangle)$,
\item le noyau de
$z - f^{N - t}x : H_*(R\langle x, z\rangle) \to H_*(R\langle x, z\rangle)$.

\end{enumerate}

Cette observation résulte d'un calcul direct\footnote{Indication : en posant
$z' = z - f^{N - t}x$, on voit que
$R\langle x, z\rangle = R\langle x, z'\rangle$ avec $\text{d}(z') = 0$
et que l'application $R\langle x, z'\rangle \to K_t$ est l'application qui annule $z'$.}, détail que nous omettons. On peut alors appliquer le lemme \ref{lemma-compute-cohomology-adjoin-variable}, partie (2), pour conclure.

\medskip\noindent
Au moyen de l'homomorphisme $K_N \to B_{N, t}$ d'algèbres différentielles graduées sur $R$ qui envoie $z$ sur $z$, nous pouvons considérer $B_{N, t} \to K_t$ comme un quasi-isomorphisme de $K_N$-modules différentiels gradués à gauche. Pour définir la flèche de l'énoncé, nous utilisons l'homomorphisme
$$
B_{N, t} = R\langle x, z, u \rangle \to K_n \otimes_R K_t,\quad
x \mapsto 1 \otimes x,\quad
z \mapsto f^{N - n}y \otimes 1,\quad
u \mapsto - f^{N - n - t}y \otimes x
$$
Cette définition a un sens dès que $N \geq n + t$. Un calcul simple montre que, dans l'un ou l'autre ordre de composition, la composée avec l'application de multiplication est l'application de transition.

\medskip\noindent

Pour $r > 1$, on écrit chacune des algèbres de Koszul comme un produit tensoriel d'algèbres de Koszul en $1$ variable, puis on applique la construction précédente. Autrement dit,
$$
K_t = R\langle x_1, \ldots, x_r\rangle =
R\langle x_1\rangle \otimes_R \ldots \otimes_R R\langle x_r\rangle
$$
où $x_i$ est de degré $1$ et $\text{d}(x_i) = f_i^t$. Lorsque $r > 1$, nous utilisons alors
$$
B_{N, t} =
R\langle x_1, z_1, u_1 \rangle
\otimes_R \ldots \otimes_R
R\langle x_r, z_r, u_r \rangle
$$
où $x_i, z_i$ sont de degré $1$ et $u_i$ est de degré $2$, avec
$\text{d}(x_i) = f_i^t$, $\text{d}(z_i) = f_i^N$,
$\text{d}(u_i) = z_i - f_i^{N - t}x_i$. L'application tensorielle $B_{N, t} \to K_t$ est un quasi-isomorphisme, car c'est un produit tensoriel de quasi-isomorphismes entre complexes de $R$-modules libres bornés inférieurement. Enfin, nous définissons l'application
$$
B_{N, t}
\to
K_n \otimes_R K_t =
R\langle y_1, \ldots, y_r\rangle \otimes_R
R\langle x_1, \ldots, x_r\rangle
$$
comme le produit tensoriel des applications construites dans le cas $r = 1$,
ou directement par les règles $x_i \mapsto 1 \otimes x_i$,
$z_i \mapsto f_i^{N - n}y_i \otimes 1$,
$u_i \mapsto - f_i^{N - n - t}y_i \otimes x_i$, ce qui a un sens dès que $N \geq n + t$. Nous omettons les détails.

\end{proof}







\begin{multicols}{2}[\section{Autres chapitres}]
\noindent
Préliminaires
\begin{enumerate}
\item \hyperref[introduction-section-phantom]{Introduction}
\item \hyperref[conventions-section-phantom]{Conventions}
\item \hyperref[sets-section-phantom]{Théorie des ensembles}
\item \hyperref[categories-section-phantom]{Catégories}
\item \hyperref[topology-section-phantom]{Topologie}
\item \hyperref[sheaves-section-phantom]{Faisceaux sur les espaces}
\item \hyperref[sites-section-phantom]{Sites et faisceaux}
\item \hyperref[stacks-section-phantom]{Champs}
\item \hyperref[fields-section-phantom]{Corps}
\item \hyperref[algebra-section-phantom]{Algèbre commutative}
\item \hyperref[brauer-section-phantom]{Groupes de Brauer}
\item \hyperref[homology-section-phantom]{Algèbre homologique}
\item \hyperref[derived-section-phantom]{Catégories dérivées}
\item \hyperref[simplicial-section-phantom]{Méthodes simpliciales}
\item \hyperref[more-algebra-section-phantom]{Compléments d'algèbre}
\item \hyperref[smoothing-section-phantom]{Lissage des morphismes d'anneaux}
\item \hyperref[modules-section-phantom]{Faisceaux de modules}
\item \hyperref[sites-modules-section-phantom]{Modules sur les sites}
\item \hyperref[injectives-section-phantom]{Objets injectifs}
\item \hyperref[cohomology-section-phantom]{Cohomologie des faisceaux}
\item \hyperref[sites-cohomology-section-phantom]{Cohomologie sur les sites}
\item \hyperref[dga-section-phantom]{Algèbre différentielle graduée}
\item \hyperref[dpa-section-phantom]{Algèbre à puissances divisées}
\item \hyperref[sdga-section-phantom]{Faisceaux différentiels gradués}
\item \hyperref[hypercovering-section-phantom]{Hyperrecouvrements}
\end{enumerate}
Schémas
\begin{enumerate}
\setcounter{enumi}{25}
\item \hyperref[schemes-section-phantom]{Schémas}
\item \hyperref[constructions-section-phantom]{Constructions de schémas}
\item \hyperref[properties-section-phantom]{Propriétés des schémas}
\item \hyperref[morphisms-section-phantom]{Morphismes de schémas}
\item \hyperref[coherent-section-phantom]{Cohomologie des schémas}
\item \hyperref[divisors-section-phantom]{Diviseurs}
\item \hyperref[limits-section-phantom]{Limites de schémas}
\item \hyperref[varieties-section-phantom]{Variétés}
\item \hyperref[topologies-section-phantom]{Topologies sur les schémas}
\item \hyperref[descent-section-phantom]{Descente}
\item \hyperref[perfect-section-phantom]{Catégories dérivées de schémas}
\item \hyperref[more-morphisms-section-phantom]{Compléments sur les morphismes}
\item \hyperref[flat-section-phantom]{Compléments sur la platitude}
\item \hyperref[groupoids-section-phantom]{Schémas en groupoïdes}
\item \hyperref[more-groupoids-section-phantom]{Compléments sur les schémas en groupoïdes}
\item \hyperref[etale-section-phantom]{Morphismes étales de schémas}
\end{enumerate}
Thèmes de théorie des schémas
\begin{enumerate}
\setcounter{enumi}{41}
\item \hyperref[chow-section-phantom]{Homologie de Chow}
\item \hyperref[intersection-section-phantom]{Théorie de l'intersection}
\item \hyperref[pic-section-phantom]{Schémas de Picard des courbes}
\item \hyperref[weil-section-phantom]{Théories cohomologiques de Weil}
\item \hyperref[adequate-section-phantom]{Modules adéquats}
\item \hyperref[dualizing-section-phantom]{Complexes dualisants}
\item \hyperref[duality-section-phantom]{Dualité pour les schémas}
\item \hyperref[discriminant-section-phantom]{Discriminants et différentes}
\item \hyperref[derham-section-phantom]{Cohomologie de de Rham}
\item \hyperref[local-cohomology-section-phantom]{Cohomologie locale}
\item \hyperref[algebraization-section-phantom]{Géométrie algébrique et formelle}
\item \hyperref[curves-section-phantom]{Courbes algébriques}
\item \hyperref[resolve-section-phantom]{Résolution des surfaces}
\item \hyperref[models-section-phantom]{Réduction semi-stable}
\item \hyperref[functors-section-phantom]{Foncteurs et morphismes}
\item \hyperref[equiv-section-phantom]{Catégories dérivées de variétés}
\item \hyperref[pione-section-phantom]{Groupes fondamentaux de schémas}
\item \hyperref[etale-cohomology-section-phantom]{Cohomologie étale}
\item \hyperref[crystalline-section-phantom]{Cohomologie cristalline}
\item \hyperref[proetale-section-phantom]{Cohomologie pro-étale}
\item \hyperref[relative-cycles-section-phantom]{Cycles relatifs}
\item \hyperref[more-etale-section-phantom]{Compléments de cohomologie étale}
\item \hyperref[trace-section-phantom]{La formule des traces}
\end{enumerate}
Espaces algébriques
\begin{enumerate}
\setcounter{enumi}{64}
\item \hyperref[spaces-section-phantom]{Espaces algébriques}
\item \hyperref[spaces-properties-section-phantom]{Propriétés des espaces algébriques}
\item \hyperref[spaces-morphisms-section-phantom]{Morphismes d'espaces algébriques}
\item \hyperref[decent-spaces-section-phantom]{Espaces algébriques décents}
\item \hyperref[spaces-cohomology-section-phantom]{Cohomologie des espaces algébriques}
\item \hyperref[spaces-limits-section-phantom]{Limites d'espaces algébriques}
\item \hyperref[spaces-divisors-section-phantom]{Diviseurs sur les espaces algébriques}
\item \hyperref[spaces-over-fields-section-phantom]{Espaces algébriques sur des corps}
\item \hyperref[spaces-topologies-section-phantom]{Topologies sur les espaces algébriques}
\item \hyperref[spaces-descent-section-phantom]{Descente et espaces algébriques}
\item \hyperref[spaces-perfect-section-phantom]{Catégories dérivées d'espaces}
\item \hyperref[spaces-more-morphisms-section-phantom]{Compléments sur les morphismes d'espaces}
\item \hyperref[spaces-flat-section-phantom]{Platitude sur les espaces algébriques}
\item \hyperref[spaces-groupoids-section-phantom]{Groupoïdes dans les espaces algébriques}
\item \hyperref[spaces-more-groupoids-section-phantom]{Compléments sur les groupoïdes dans les espaces}
\item \hyperref[bootstrap-section-phantom]{Amorçage}
\item \hyperref[spaces-pushouts-section-phantom]{Sommes amalgamées d'espaces algébriques}
\end{enumerate}
Thèmes de géométrie
\begin{enumerate}
\setcounter{enumi}{81}
\item \hyperref[spaces-chow-section-phantom]{Groupes de Chow des espaces}
\item \hyperref[groupoids-quotients-section-phantom]{Quotients de groupoïdes}
\item \hyperref[spaces-more-cohomology-section-phantom]{Compléments sur la cohomologie des espaces}
\item \hyperref[spaces-simplicial-section-phantom]{Espaces simpliciaux}
\item \hyperref[spaces-duality-section-phantom]{Dualité pour les espaces}
\item \hyperref[formal-spaces-section-phantom]{Espaces algébriques formels}
\item \hyperref[restricted-section-phantom]{Algébrisation des espaces formels}
\item \hyperref[spaces-resolve-section-phantom]{Résolution des surfaces revisitée}
\end{enumerate}
Théorie des déformations
\begin{enumerate}
\setcounter{enumi}{89}
\item \hyperref[formal-defos-section-phantom]{Théorie formelle des déformations}
\item \hyperref[defos-section-phantom]{Théorie des déformations}
\item \hyperref[cotangent-section-phantom]{Le complexe cotangent}
\item \hyperref[examples-defos-section-phantom]{Problèmes de déformation}
\end{enumerate}
Champs algébriques
\begin{enumerate}
\setcounter{enumi}{93}
\item \hyperref[algebraic-section-phantom]{Champs algébriques}
\item \hyperref[examples-stacks-section-phantom]{Exemples de champs}
\item \hyperref[stacks-sheaves-section-phantom]{Faisceaux sur les champs algébriques}
\item \hyperref[criteria-section-phantom]{Critères de représentabilité}
\item \hyperref[artin-section-phantom]{Axiomes d'Artin}
\item \hyperref[quot-section-phantom]{Espaces de Quot et de Hilbert}
\item \hyperref[stacks-properties-section-phantom]{Propriétés des champs algébriques}
\item \hyperref[stacks-morphisms-section-phantom]{Morphismes de champs algébriques}
\item \hyperref[stacks-limits-section-phantom]{Limites de champs algébriques}
\item \hyperref[stacks-cohomology-section-phantom]{Cohomologie des champs algébriques}
\item \hyperref[stacks-perfect-section-phantom]{Catégories dérivées de champs}
\item \hyperref[stacks-introduction-section-phantom]{Introduction aux champs algébriques}
\item \hyperref[stacks-more-morphisms-section-phantom]{Compléments sur les morphismes de champs}
\item \hyperref[stacks-geometry-section-phantom]{La géométrie des champs}
\end{enumerate}
Thèmes de théorie des espaces de modules
\begin{enumerate}
\setcounter{enumi}{107}
\item \hyperref[moduli-section-phantom]{Champs de modules}
\item \hyperref[moduli-curves-section-phantom]{Espaces de modules de courbes}
\end{enumerate}
Divers
\begin{enumerate}
\setcounter{enumi}{109}
\item \hyperref[examples-section-phantom]{Exemples}
\item \hyperref[exercises-section-phantom]{Exercices}
\item \hyperref[guide-section-phantom]{Guide bibliographique}
\item \hyperref[desirables-section-phantom]{Desiderata}
\item \hyperref[coding-section-phantom]{Style de programmation}
\item \hyperref[obsolete-section-phantom]{Obsolète}
\item \hyperref[fdl-section-phantom]{Licence de documentation libre GNU}
\item \hyperref[index-section-phantom]{Index généré automatiquement}
\end{enumerate}
\end{multicols}


\bibliography{my} \bibliographystyle{amsalpha}

\end{document}
